
\documentclass[11pt]{article}

\usepackage[a4paper,margin=1in]{geometry}
\usepackage{fontspec}
\usepackage{microtype}
\usepackage{amsmath,amssymb,amsthm,mathtools}
\usepackage{booktabs,longtable,array}
\usepackage{enumitem}
\usepackage{xurl}
\usepackage[hidelinks]{hyperref}
\usepackage{bookmark}

\setlength{\parskip}{0.45em}
\setlength{\parindent}{0pt}
\setcounter{secnumdepth}{3}
\setcounter{tocdepth}{2}
\setlength{\LTleft}{0pt}
\setlength{\LTright}{0pt}
\sloppy
\providecommand{\tightlist}{%
  \setlength{\itemsep}{0pt}\setlength{\parskip}{0pt}}

\theoremstyle{definition}
\newtheorem{theorem}{Theorem}[section]
\newtheorem{lemma}[theorem]{Lemma}
\newtheorem{proposition}[theorem]{Proposition}
\newtheorem{corollary}[theorem]{Corollary}
\renewcommand{\qedsymbol}{$\square$}

\makeatletter
\renewcommand{\verbatim@font}{\normalfont\ttfamily\small}
\makeatother


\title{\(\sigma\)-Equivariant magmas on \(\operatorname{AG}(2,3)\): finite landscape, intrinsic selection, and linearization}
\author{Burundai Taryi\thanks{Independent researcher, \href{mailto:burundai.taryi@gmail.com}{burundai.taryi@gmail.com}.}}
\date{}

\begin{document}
\maketitle

\begin{abstract}
We study a finite family of \(\sigma\)-equivariant magmas on \[
M=S\times S,\qquad S=\mathbb F_3,
\] with left-row output and a fixed same-row Steiner rule. The normal
form is \[
\Omega'\cong S^{18}\times S^3\cong S^{21},\qquad |\Omega'|=3^{21}.
\]

The paper has three layers. First, it maps the finite landscape: the
certified associativity atlas gives \[
63\le \operatorname{Assoc}\le 597,
\] with endpoint loci and exact partition function, while the global
hidden-continuation theorem gives \[
\min H_{tot}=-2268,\qquad \max H_{tot}=7302,
\qquad \max_{N_-=0}H_{tot}=7020.
\] The layer above the pure frontier consists of twelve impure
signed-cancellation maximizers.

Second, the central selection theorem isolates one magma by four
internal finite conditions: \[
\boxed{
\mathcal H_{acc}=0,
\qquad
H_{diag}=0,
\qquad
\text{diagonal idempotence},
\qquad
\text{pure }C/J.
}
\] These conditions force the PAB multiplication \[
(r_1,c_1)\cdot(r_2,c_2)=
\begin{cases}
(r_1,r_2),&r_1\ne r_2,\\[1mm]
(r_1,\overline{c_1c_2}),&r_1=r_2,
\ c_1\ne c_2,\\[1mm]
(r_1,r_1),&r_1=r_2,
\ c_1=c_2.
\end{cases}
\]

Third, after selection, PAB admits a linear/operator passport:
rank-three left translations, a finite composition envelope,
endomorphism and intertwiner structure, a binary term algebra of size \[
|T_2(\mathrm{PAB})|=630,
\] a row-contraction matrix bridge, a hidden-continuation operator
identity, and companion carrier-level Lie and signed symplectic
structures.

The organizing implication is \[
\boxed{\text{finite landscape}}
\Longrightarrow
\boxed{\text{intrinsic finite selection}}
\Longrightarrow
\boxed{\text{linearization of the selected magma}}.
\]
\end{abstract}

\tableofcontents
\newpage

\setcounter{section}{-1}
\section{Introduction}\label{introduction}

\subsection{Aim}\label{aim}

The starting point is a finite landscape of multiplications on \[
M=S\times S,\qquad S=\mathbb F_3.
\] The first coordinate is a row, the second is a column. The family is
defined by left-row output, the same-row Steiner rule, and
\(\sigma\)-equivariance. These axioms leave exactly twenty-one ternary
coordinates, \[
\Omega'\cong S^{18}\times S^3\cong S^{21},
\qquad |\Omega'|=3^{21}=10\,460\,353\,203.
\]

The manuscript follows one spine: \[
\text{finite landscape}
\longrightarrow
\text{intrinsic finite selection}
\longrightarrow
\text{linearization after selection}.
\] The global atlases describe the ambient landscape; the compact
selector isolates PAB; the operator, matrix-carrier, Lie, symplectic,
and term-algebra results describe the selected magma.

\subsection{Main result: compact
selection}\label{main-result-compact-selection}

The central theorem is the uniqueness of PAB under four finite
conditions: \[
\boxed{
\mathcal H_{acc}=0,
\qquad
H_{diag}=0,
\qquad
\text{diagonal idempotence},
\qquad
\text{pure }C/J.
}
\] Here \[
\mathcal H_{acc}=I(C_{out};C_1,C_2\mid R_1,R_2)
\] is zero hidden-column access in the cross-row channel, and \[
H_{diag}=-\sum_s p_s\log_2p_s
\] is the Shannon entropy of the same-point diagonal output multiset.
Their zero sets are the finite structural conditions of column-blindness
and diagonal constancy. Diagonal idempotence fixes the constant diagonal
to \(d=000\). The final pure \(C/J\) condition is imposed on the six
directed edges \[
M^\times=\{(r,c):r\ne c\}.
\]

The selected multiplication is \[
(r_1,c_1)\cdot(r_2,c_2)=
\begin{cases}
(r_1,r_2), & r_1\ne r_2,\\[2mm]
(r_1,\overline{c_1c_2}), & r_1=r_2,
\ c_1\ne c_2,\\[2mm]
(r_1,r_1), & r_1=r_2,
\ c_1=c_2.
\end{cases}
\] A longer explanatory chain is also retained, but it is not part of the
compact logical core: \[
\Omega'
\xrightarrow{\mathcal H_{acc}=0,
\ H_{diag}=0}
9\times3
\xrightarrow{a\ne b,
\ d\text{ idempotent}}
6\times1
\xrightarrow{\min\operatorname{Assoc}_{000}}
2\times1
\xrightarrow{\text{pure }C/J}
1.
\] It displays the PAB/row-complement ambiguity before the undrifted
\(C/J\) separation; Theorem~\ref{thm:12.3} gives the shorter selection
proof in which pure \(C/J\), after the preceding three selectors, already
forces PAB.

\subsection{Landscape atlases}\label{landscape-atlases}

The associativity statistic \[
\operatorname{Assoc}(x)=
\#\{(u,v,w):(u\cdot v)\cdot w=u\cdot(v\cdot w)\}
\] has global range \[
63\le\operatorname{Assoc}\le597.
\] Its exact partition function has \(167\) nonzero coefficients, mean
\(245\), variance \(27820/27\), and mode \(237\) with coefficient
\(426,317,976\). PAB has \[
\operatorname{Assoc}(\mathrm{PAB})=219.
\]

The hidden-continuation atlas uses local signed contributions \(h_q\),
their total \[
H_{tot}=3\sum_{q\in S^5}h_q,
\] and the negative local count \(N_-\). The global theorem gives \[
\min_{\Omega'}H_{tot}=-2268,
\qquad
\max_{\Omega'}H_{tot}=7302,
\qquad
\max_{N_-=0}H_{tot}=7020.
\] PAB and row-complement lie on this pure frontier; the twelve points
above it are impure signed-cancellation maximizers.

\subsection{Linearization after
selection}\label{linearization-after-selection}

After selection, PAB has a compact operator theory. Every left
translation has image a row fiber and rank \(3\). The positive
composition envelope has \(51\) maps, decomposed into three row-target
components of size \(17\), and its linear span has dimension \(27\). The
endomorphism monoid consists of the six simultaneous symbol permutations
and three diagonal collapses. The binary term algebra has normal form \[
T_2(\mathrm{PAB})
\cong
\{x,y\}\times\mathcal S_{21}\times\mathcal U_{15},
\qquad |T_2(\mathrm{PAB})|=630.
\]

The vector-space carrier identifies with matrix units, \[
k[M]\cong\operatorname{Mat}_3(k),
\qquad e_{(r,c)}\leftrightarrow E_{rc}.
\] This carrier supports the row-contraction bridge \[
\mu_{cross}(A,B)=AJ_3B^\top=(A\mathbf1)(B\mathbf1)^\top,
\] the compact carrier-level type \(SU(3)\times SU(2)\times U(1)\), and
a signed symplectic lift of the pure \(C/J\) directed-edge geometry.

\subsection{Proof architecture and
support}\label{proof-architecture-and-support}

The proof has three support modes.

\begin{longtable}[]{@{}
  >{\raggedright\arraybackslash}p{(\columnwidth - 4\tabcolsep) * \real{0.3333}}
  >{\raggedright\arraybackslash}p{(\columnwidth - 4\tabcolsep) * \real{0.3333}}
  >{\raggedright\arraybackslash}p{(\columnwidth - 4\tabcolsep) * \real{0.3333}}@{}}
\toprule\noalign{}
\begin{minipage}[b]{\linewidth}\raggedright
result family
\end{minipage} & \begin{minipage}[b]{\linewidth}\raggedright
role
\end{minipage} & \begin{minipage}[b]{\linewidth}\raggedright
support
\end{minipage} \\
\midrule\noalign{}
\endhead
\bottomrule\noalign{}
\endlastfoot
normal form, access-zero, diagonal compression, compact PAB selector &
hand finite theorems & text; L2 tables audit survivor counts \\
associativity range, endpoints, partition function & finite landscape
atlas & compact \texttt{tables/} plus L1 verifier \\
global hidden-continuation frontier & global frontier certificate &
separate \texttt{mathcal\_H/} bundle \\
operator, endomorphism, composition, \(T_2\) structure & post-selection
structure & text plus L3 realization/transition/audit tables \\
carrier Lie and signed symplectic structures & downstream bridge facts &
bridge audit tables \\
\end{longtable}

The distributed support corpus keeps these roles separate: \[
\texttt{scripts/},\qquad \texttt{tables/},\qquad \texttt{mathcal\_H/}.
\] The compact manuscript-facing check is

\begin{verbatim}
python3 -S scripts/verify_preprint_support.py
\end{verbatim}

The global hidden-continuation bundle is checked separately from inside
\texttt{mathcal\_H/}.

\subsection{Positioning}\label{positioning}

The first selection layer uses standard finite information observables
in the language of Landauer, Bennett, Shannon, and Cover--Thomas
\cite{Landauer1961,Bennett1973,Bennett1982,Shannon1948,CoverThomas2006}.
The finite magma, term-function, and clone-theoretic context is standard
universal algebra
\cite{Birkhoff1935,Cohn1965,BurrisSankappanavar1981,Post1941,Rosenberg1970,Szendrei1986,Lau2006}.
The same-row rule is the three-point Steiner multiplication, in the
classical nonassociative and quasigroup background
\cite{Bruck1958,DenesKeedwell1974,Pflugfelder1990,ColbournDinitz2007}.
The certificate layer belongs to the tradition of computer-assisted
finite mathematics
\cite{AppelHaken1977,AppelHakenKoch1977,Gonthier2008,McCune1997,Hales2005,HalesEtAl2017}.
The carrier-level matrix, Lie, and symplectic terminology follows
standard conventions
\cite{HornJohnson2012,FultonHarris1991,Hall2015,GuilleminSternberg1990,McDuffSalamon2017}.

\subsection{Notation}\label{notation}

\begin{longtable}[]{@{}
  >{\raggedright\arraybackslash}p{(\columnwidth - 2\tabcolsep) * \real{0.5000}}
  >{\raggedright\arraybackslash}p{(\columnwidth - 2\tabcolsep) * \real{0.5000}}@{}}
\toprule\noalign{}
\begin{minipage}[b]{\linewidth}\raggedright
notation
\end{minipage} & \begin{minipage}[b]{\linewidth}\raggedright
meaning
\end{minipage} \\
\midrule\noalign{}
\endhead
\bottomrule\noalign{}
\endlastfoot
\(S\) & the field \(\mathbb F_3\) \\
\(M\) & the carrier \(S\times S\) \\
\(F_r\) & the row fiber \(\{r\}\times S\) \\
\(M^\times\) & the six off-diagonal points \(\{(r,c):r\ne c\}\) \\
\(\Omega'\) & the AX0+AX2, \(\sigma\)-equivariant landscape \\
\((A,B,d)\) & normal coordinates in \(S^{18}\times S^3\) \\
\(h_{a,b}\) & column-blind cross-row rule with anchors \(h(0,1)=a\),
\(h(0,2)=b\) \\
\(\operatorname{Assoc}\) & total number of associative triples in
\(M^3\) \\
\(\operatorname{Assoc}_{000}\) & local associativity statistic after
fixing \(d=000\) \\
\(\mathcal H_{acc}\) & hidden-column access observable \\
\(H_{diag}\) & entropy of the diagonal output multiset \\
\(h_q,H_{tot},N_-\) & hidden-continuation local contribution, total
value, and negative local count \\
\(C,J\) & continuation and reversal maps on \(M^\times\) \\
\(T_2(\mathrm{PAB})\) & binary term-function algebra of PAB \\
\end{longtable}

\part{The finite landscape}

\section{Carrier, axioms, and normal
form}\label{carrier-axioms-and-normal-form}

\subsection{The row-graded carrier}\label{the-row-graded-carrier}

Let \[
S=\mathbb F_3,
\qquad
M=S\times S.
\] The first coordinate is called the row and the second coordinate the
column. The row fibers are \(F_r=\{r\}\times S\).

The left-row output axiom says that the product of \((r_1,c_1)\) and
\((r_2,c_2)\) always lies in \(F_{r_1}\). Hence a multiplication in the
family is determined by its output-column function.

The same-row off-diagonal axiom fixes the product inside a row when the
two columns are distinct: \[
(r,i)\cdot(r,j)=(r,\overline{ij})
\qquad (i\ne j).
\] Here \(\overline{ij}\) denotes the unique element of
\(S\setminus\{i,j\}\). This is the Steiner multiplication on a
three-point row fiber \cite{Pflugfelder1990,ColbournDinitz2007}.

\subsection{Normal coordinates}\label{normal-coordinates}

The remaining freedom has two kinds. First, cross-row products are free
in their output column. By \(\sigma\)-equivariance, it is enough to
record two representative cross-row tables, denoted \(A\) and \(B\),
each with nine entries. Second, same-point products in a row are
governed by a diagonal vector \[
d=(d_0,d_1,d_2)\in S^3.
\] The normal coordinates are therefore \[
(A,B,d)\in S^9\times S^9\times S^3.
\]

\subsection{Normal form theorem}\label{normal-form-theorem}

\begin{theorem}
\label{thm:1.1}
The active AX0+AX2, \(\sigma\)-equivariant
landscape has normal form \[
\Omega'\cong S^{18}\times S^3\cong S^{21}.
\] Consequently \[
|\Omega'|=3^{21}=10\,460\,353\,203.
\]
\end{theorem}

\begin{proof}
AX0 fixes the output row of every product. AX2 fixes all
same-row off-diagonal output columns. For cross-row products,
\(\sigma\)-equivariance reduces the row-pair data to the two ordered
representatives \((0,1)\) and \((0,2)\). Each representative has
\(3\times3\) column-input entries, giving \(18\) cross-row coordinates.
The same-point diagonal sector contributes \(3\) coordinates. These
coordinates are independent, so the landscape is \(S^{21}\).
\end{proof}

\section{Effective symmetry and orbit
count}\label{effective-symmetry-and-orbit-count}

The resulting family has a residual effective symbol-relabelling
symmetry. Because the landscape is already written in
\(\sigma\)-equivariant normal coordinates, the effective action on these
coordinates is generated by the remaining nontrivial relabelling class.
We denote this residual action by \[
G_{\mathrm{eff}}\cong C_2.
\] The diagonal sector still carries the ordinary symbol-relabelling
compression \[
|\Delta/\operatorname{Sym}(S)|=15.
\]

\begin{proposition}
\label{prop:2.1}
The effective orbit count of the full
normal-coordinate landscape is \[
|\Omega'/G_{\mathrm{eff}}|=\frac{3^{21}+3^{10}}2.
\]
\end{proposition}

\begin{proof}
The computation is Burnside's lemma for the residual
effective action on normal coordinates. The identity contributes
\(3^{21}\). The nontrivial effective symmetry has fixed locus of size
\(3^{10}\). Dividing by \(|G_{\mathrm{eff}}|=2\) gives the formula.
\end{proof}

The orbit count is a landscape compression.

\section{Hidden-column access}\label{hidden-column-access}

The first selection functional is information-theoretic. It measures
whether the cross-row rule has functional dependence on hidden column
coordinates after the row pair is known. Its vanishing is the
Landauer-facing condition that the cross-row channel admits a
realization without querying column input.

\subsection{Column-blind cross-row
rules}\label{column-blind-cross-row-rules}

A cross-row output rule is \textbf{column-blind} if its output column
depends only on the two row coordinates: \[
g(r_1,c_1,r_2,c_2)=h(r_1,r_2).
\] Equivalently, after the ordered row pair is fixed, the output column
is constant over the nine column inputs.

By \(\sigma\)-equivariance, a column-blind rule is determined by two
values \[
a=h(0,1),
\qquad
b=h(0,2).
\] We call \((a,b)\) the cross-row anchors and write \(h_{a,b}\) for the
corresponding rule. The remaining ordered row pairs are forced by
\(\sigma\)-shift: \[
h(r,r+1)=r+a,
\qquad
h(r,r+2)=r+b
\] with arithmetic in \(\mathbb F_3\). There are \(3^2=9\) column-blind
cross-row rules.

\subsection{Hidden-column access as a Landauer-facing
observable}\label{hidden-column-access-as-a-landauer-facing-observable}

Let \[
\Sigma_{cross}=\{(r_1,c_1,r_2,c_2)\in S^4:r_1\ne r_2\}
\] with the uniform distribution, and let \[
C_{out}=g(R_1,C_1,R_2,C_2)
\] be the output-column random variable. Define the hidden-column access
observable by \[
\boxed{
\mathcal H_{acc}(g)=I(C_{out};C_1,C_2\mid R_1,R_2).
}
\] Logarithms are base \(2\), so the value is measured in bits.

This is a structural information observable. It measures how much
hidden-column information is functionally retained in the cross-row
output once the row pair is fixed. The Landauer reading is that hidden
column information which is read and later discarded by an irreversible
implementation carries thermodynamic cost
\cite{Landauer1961,Bennett1973,Bennett1982}. Therefore
\(\mathcal H_{acc}=0\) selects exactly those cross-row rules for which
that cost is structurally avoidable: the output can be computed without
consulting the hidden columns.

The zero condition is used in this implementation-independent structural
sense: the column-read burden can be avoided altogether.

\subsection{Access-zero theorem}\label{access-zero-theorem}

\begin{theorem}
\label{thm:3.1}
For a cross-row rule \(g\), \[
\mathcal H_{acc}(g)=0
\quad\Longleftrightarrow\quad
 g(r_1,c_1,r_2,c_2)=h(r_1,r_2).
\] There are exactly nine column-blind cross-row rules.
\end{theorem}

\begin{proof}
Conditional mutual information vanishes exactly when \[
C_{out}\perp (C_1,C_2)\mid (R_1,R_2).
\] Since \(C_{out}\) is a deterministic function of
\((R_1,C_1,R_2,C_2)\), conditional independence is equivalent to
constancy in \((C_1,C_2)\) for every fixed ordered row pair
\((R_1,R_2)\). That is precisely column-blindness.

The two representative cross-row row pairs \((0,1)\) and \((0,2)\) each
receive one output-column value in \(S\), giving \(3^2=9\)
\(\sigma\)-equivariant column-blind rules.
\end{proof}

This theorem is a landscape result and also the first reduction in the
selection chain.

\section{Associativity master
formula}\label{associativity-master-formula}

\subsection{The associativity
statistic}\label{the-associativity-statistic}

For \(x\in\Omega'\), define \[
\operatorname{Assoc}(x)=
\#\{(u,v,w)\in M^3:(u\cdot v)\cdot w=u\cdot(v\cdot w)\}.
\] Since \(|M|=9\), this statistic lies between \(0\) and \(729\).

\begin{theorem}
\label{thm:4.1}
In normal coordinates, \[
\operatorname{Assoc}=3\,rawAssoc,
\] where \(rawAssoc\) is a normalized sum over \(3^5=243\) triples and
decomposes as \[
rawAssoc=rawRRR+rawRRS+rawRSR+rawSRR+rawDIST.
\] The five blocks correspond to the row patterns \[
RRR,
\qquad RRS,
\qquad RSR,
\qquad SRR,
\qquad DIST.
\]
\end{theorem}

\begin{proof}
By \(\sigma\)-equivariance the first row coordinate of a
triple may be normalized to \(0\). The two remaining row variables and
the three column variables give \(3^5\) normalized cases; the factor
\(3\) restores the three translates. The displayed blocks are the five
equality patterns of the three rows.
\end{proof}

\subsection{The same-row block}\label{the-same-row-block}

The all-same-row block depends only on the diagonal vector \(d\). Let \[
E(d)=\#\{(i,j):i\ne j,
\ d_i=d_j\},
\qquad
N(d)=\#\{(i,j):i\ne j,
\ d_i=j,
\ d_j=j\}.
\]

\begin{lemma}
\label{lem:4.2}
The same-row contribution is \[
rawRRR(d)=9+E(d)+2N(d).
\] The four diagonal classes give raw same-row values \[
9,\qquad 11,\qquad 13,\qquad 19.
\] In particular, \[
rawRRR(000)=19.
\]
\end{lemma}

\begin{proof}
In one row, unequal column pairs use the Steiner rule
and equal pairs use \(d\). The \(27\) triples split into a baseline
family plus the two diagonal-sensitive families counted by \(E(d)\) and
\(N(d)\). This gives the formula and the four class values.
\end{proof}

\subsection{\texorpdfstring{Column-blind profile at
\(d=000\)}{Column-blind profile at d=000}}\label{column-blind-profile-at-d000}

Let \(h_{a,b}\) be the column-blind rule \[
h(r,r+1)=r+a,
\qquad
h(r,r+2)=r+b.
\] At \(d=000\), same-point products satisfy \((r,c)^2=(r,r)\).

\begin{proposition}
\label{prop:4.3}
On the column-blind, \(d=000\) stratum, \[
\begin{aligned}
&(rawRRR,rawRRS,rawRSR,rawSRR,rawDIST)(a,b,000)\\
&\quad =
\bigl(19,\, 9[a=0]+9[b=0],\, 9[a=0]+9[b=0],\\
&\hspace{7em}54,\,54[a=b]\bigr).
\end{aligned}
\] Consequently \[
\boxed{
\operatorname{Assoc}_{000}(h_{a,b})
=
3\left(73+18[a=0]+18[b=0]+54[a=b]\right).
}
\]
\end{proposition}

\begin{proof}
The \(RRR\) value is Lemma 4.2. In the \(RRS\) and
\(RSR\) patterns, the same-row step must return the cross-row anchor; at
\(d=000\) this contributes \(9\) exactly when the relevant anchor is
\(0\). In \(SRR\), the final cross-row output ignores same-row column
processing, giving all \(54\) normalized triples. In \(DIST\), the two
association orders compare the two cross-row anchors, so all \(54\)
triples survive exactly when \(a=b\).
\end{proof}

The nine column-blind rules split as follows:

\begin{longtable}[]{@{}
  >{\raggedright\arraybackslash}p{(\columnwidth - 6\tabcolsep) * \real{0.2400}}
  >{\raggedright\arraybackslash}p{(\columnwidth - 6\tabcolsep) * \real{0.3000}}
  >{\raggedleft\arraybackslash}p{(\columnwidth - 6\tabcolsep) * \real{0.2600}}
  >{\raggedleft\arraybackslash}p{(\columnwidth - 6\tabcolsep) * \real{0.2000}}@{}}
\toprule\noalign{}
\begin{minipage}[b]{\linewidth}\raggedright
class
\end{minipage} & \begin{minipage}[b]{\linewidth}\raggedright
representatives
\end{minipage} & \begin{minipage}[b]{\linewidth}\raggedleft
profile
\end{minipage} & \begin{minipage}[b]{\linewidth}\raggedleft
\(\operatorname{Assoc}_{000}\)
\end{minipage} \\
\midrule\noalign{}
\endhead
\bottomrule\noalign{}
\endlastfoot
anchored trivial & \((0,0)\) & \((19,18,18,54,54)\) & \(489\) \\
off-anchor trivial & \((1,1),(2,2)\) & \((19,0,0,54,54)\) & \(381\) \\
partly anchored nondegenerate & \((0,1),(0,2),(1,0),(2,0)\) &
\((19,9,9,54,0)\) & \(273\) \\
fully off-anchor nondegenerate & \((1,2),(2,1)\) & \((19,0,0,54,0)\) &
\(219\) \\
\end{longtable}

The fourth class consists of PAB and row-complement.

\section{Global associativity range}\label{global-associativity-range}

\subsection{Exact range}\label{exact-range}

\begin{lemma}[Associativity atlas certificate protocol]
\label{lem:5.1}
The
global associativity atlas is used through five finite obligations.

\begin{enumerate}
\def\labelenumi{\arabic{enumi}.}
\tightlist
\item
  \textbf{A1, block reduction.} The associativity statistic is reduced
  by \(\sigma\)-normalization to the five row-pattern blocks \[
  rawAssoc=rawRRR+rawRRS+rawRSR+rawSRR+rawDIST
  \] of Theorem 4.1.
\item
  \textbf{A2, diagonal decomposition.} The same-row block is evaluated
  from the diagonal vector by Lemma 4.2, giving the four possible values
  \[
  9,\quad 11,\quad 13,\quad 19.
  \]
\item
  \textbf{A3, fixed-diagonal cross-range certificates.} For each
  diagonal class, the certificate verifies the exact range of \[
  rawCross=rawRRS+rawRSR+rawSRR+rawDIST.
  \]
\item
  \textbf{A4, endpoint witnesses and endpoint counts.} The endpoint
  certificates verify the normal-coordinate witnesses, endpoint counts,
  and orbit representatives for \(\operatorname{Assoc}=63\) and
  \(\operatorname{Assoc}=597\).
\item
  \textbf{A5, distribution consistency.} The partition-function
  certificate gives the full coefficient table for \[
  Z(q)=\sum_{x\in\Omega'}q^{\operatorname{Assoc}(x)}
  \] and checks total mass, range, moments, and mode.
\end{enumerate}

Any certificate satisfying A1--A5 proves the global associativity range,
endpoint loci, and partition-function claims of Sections 5--6. This is a
certified atlas statement for the ambient landscape; the compact PAB
uniqueness proof in Theorem 12.3 is the separate hand spine in Part II.
\hfill$\square$
\end{lemma}

\begin{theorem}[Global associativity atlas]
\label{thm:5.2}
On \(\Omega'\), \[
21\le rawAssoc\le199.
\] Equivalently, \[
63\le\operatorname{Assoc}\le597.
\]

The finite certificate is organized by diagonal class:

\begin{longtable}[]{@{}
  >{\raggedright\arraybackslash}p{(\columnwidth - 6\tabcolsep) * \real{0.2000}}
  >{\raggedleft\arraybackslash}p{(\columnwidth - 6\tabcolsep) * \real{0.2667}}
  >{\raggedleft\arraybackslash}p{(\columnwidth - 6\tabcolsep) * \real{0.2667}}
  >{\raggedleft\arraybackslash}p{(\columnwidth - 6\tabcolsep) * \real{0.2667}}@{}}
\toprule\noalign{}
\begin{minipage}[b]{\linewidth}\raggedright
diagonal class
\end{minipage} & \begin{minipage}[b]{\linewidth}\raggedleft
rawRRR
\end{minipage} & \begin{minipage}[b]{\linewidth}\raggedleft
certified rawCross range
\end{minipage} & \begin{minipage}[b]{\linewidth}\raggedleft
rawAssoc consequence
\end{minipage} \\
\midrule\noalign{}
\endhead
\bottomrule\noalign{}
\endlastfoot
permutation & 9 & 12..180 & 21..189 \\
two-valued, no fixed point & 11 & 10..162 & 21..173 \\
two-valued, with fixed point & 13 & 12..174 & 25..187 \\
constant & 19 & 12..180 & 31..199 \\
\end{longtable}

Here \[
rawCross=rawRRS+rawRSR+rawSRR+rawDIST.
\] The endpoint values are therefore \(rawAssoc=21\) and
\(rawAssoc=199\), or \(\operatorname{Assoc}=63\) and
\(\operatorname{Assoc}=597\).
\end{theorem}

\begin{proof}
The hand-readable reduction is Theorem 4.1 plus Lemma
4.2. After that reduction, the global claim is exactly the A3
cross-range obligation of Lemma 5.1. Adding the certified cross-row
ranges to the four possible same-row values gives the displayed global
range.

The endpoint witnesses, endpoint counts, orbit representatives, and
partition-function consistency are the A4--A5 obligations of Lemma 5.1.
They are recorded in \texttt{l1\_assoc\_distribution.csv},
\texttt{l1\_extremal\_loci\_counts.csv},
\texttt{l1\_extremal\_loci\_orbits.csv},
\texttt{l1\_extremal\_loci\_solutions.csv}, and
\texttt{l1\_assoc\_proof\_targets.csv}, and checked by the compact L1
verifier.
\end{proof}

\subsection{Global and local associativity
roles}\label{global-and-local-associativity-roles}

PAB has \[
\operatorname{Assoc}(\mathrm{PAB})=219,
\] while the landscape minimum is \(63\). The global statistic maps the
ambient landscape. The selection argument later uses the conditional
local statistic \(\operatorname{Assoc}_{000}\), after information
compression and anchoring, to expose the final two-candidate ambiguity.

\section{Endpoint loci and partition
function}\label{endpoint-loci-and-partition-function}

\subsection{Endpoint loci}\label{endpoint-loci}

\begin{theorem}
\label{thm:6.1}
The associativity endpoint loci have sizes \[
|\mathcal M_{63}|=24,
\qquad
|\mathcal M_{63}/\operatorname{Sym}(S)|=12,
\] and \[
|\mathcal M_{597}|=6,
\qquad
|\mathcal M_{597}/\operatorname{Sym}(S)|=3.
\] The minimum endpoint block profiles are \[
(11,2,2,2,4)
\quad\text{and}\quad
(9,2,2,4,4),
\] while the maximum endpoint block profile is \[
(19,36,36,54,54).
\]
\end{theorem}

\begin{proof}
This is the A4 endpoint part of Lemma 5.1. The endpoint
certificates record the exact solution counts by diagonal sector, orbit
representatives, and explicit witnesses in normal coordinates. The
minimum occurs in the permutation and two-valued/no-fixed diagonal
classes. The maximum occurs in the constant diagonal class. The block
profiles are obtained by evaluating the five normalized associativity
blocks on the endpoint witnesses.
\end{proof}

\subsection{Exact associativity partition
function}\label{exact-associativity-partition-function}

\begin{theorem}
\label{thm:6.2}
The exact polynomial \[
Z(q)=\sum_{x\in\Omega'}q^{\operatorname{Assoc}(x)}
\] has \(167\) nonzero coefficients and total mass \[
Z(1)=3^{21}.
\] Its first two moments are \[
\mathbb E[\operatorname{Assoc}]=245,
\qquad
\operatorname{Var}(\operatorname{Assoc})=\frac{27820}{27}.
\] The mode occurs at \[
\operatorname{Assoc}=237
\] with coefficient \[
426,317,976.
\]
\end{theorem}

\begin{proof}
This is the A5 partition-function part of Lemma 5.1. The
coefficient certificate \texttt{l1\_assoc\_distribution.csv} is the
exact finite distribution of \(\operatorname{Assoc}\) on \(\Omega'\).
The total mass, range, mean, variance, and mode are computed directly
from that table and checked by the compact L1 verifier.
\end{proof}

\section{Hidden continuation}\label{hidden-continuation}

\subsection{Local distance reduction}\label{local-distance-reduction}

The hidden-continuation observable is a finite path-information defect.
The continuation side keeps the two-step route data, while the
left-translation side compares only the collapsed endpoint operators.
Thus \(h_q/2\) is the signed Hamming contrast between a path-resolved
comparison and an endpoint-resolved comparison. Negative local values
are precisely the local cancellation failures counted by \(N_-\). The
normalization \(H_{tot}=6rawH\) is the \(\sigma\)-orbit normalization
used by the hidden-continuation certificate.

Its local index is
\[
q=(b,t,a,e,f)\in S^5.
\]
For a completed normal-form point \(x\in\Omega'\), let \(M_s\) denote
the local table read: \(M_1\) and \(M_2\) are the two cross-row tables,
and \(M_0\) is the same-row diagonal/complement rule. All row and column
arithmetic below is in \(S=\mathbb F_3\). Define
\[
xy=M_b(a,e),
\qquad
yz=M_{t-b}(e-b,f-b),
\qquad
z_R=b+yz,
\]
and
\[
ep_L=M_t(xy,f),
\qquad
ep_R=M_b(a,z_R).
\]

For \(c\in S\), define the collapsed left-translation signature
\[
L_c(i,j)=M_i(c,j)
\qquad ((i,j)\in S^2).
\]
For \(s,\alpha,z\in S\), define the stepwise continuation signature
\[
C_{s,\alpha,z}(i,j)
=
M_s\!\left(
\alpha,\,
s+M_{i-s}(z-s,j-s)
\right)
\qquad ((i,j)\in S^2).
\]
The associated Hamming distances are
\[
D_L(c,c')
=
\#\{(i,j)\in S^2:L_c(i,j)\ne L_{c'}(i,j)\},
\]
and
\[
D_C(s,\alpha,z;s',\alpha',z')
=
\#\{(i,j)\in S^2:
C_{s,\alpha,z}(i,j)\ne C_{s',\alpha',z'}(i,j)\}.
\]

\paragraph{Direct local definition.}
For each of the nine test pairs \((i,j)\in S^2\), set
\[
\delta^C_{ij}(q;x)
=
\mathbf 1\!\left\{
C_{t,xy,f}(i,j)\ne C_{b,a,z_R}(i,j)
\right\},
\]
and
\[
\delta^L_{ij}(q;x)
=
\mathbf 1\!\left\{
L_{ep_L}(i,j)\ne L_{ep_R}(i,j)
\right\}.
\]
The nine direct local summands are
\[
\Delta_{ij}(q;x)=2\bigl(\delta^C_{ij}(q;x)-\delta^L_{ij}(q;x)\bigr),
\]
displayed as
\[
\begin{array}{c|ccc}
 & j=0 & j=1 & j=2\\
\hline
i=0 & \Delta_{00} & \Delta_{01} & \Delta_{02}\\
i=1 & \Delta_{10} & \Delta_{11} & \Delta_{12}\\
i=2 & \Delta_{20} & \Delta_{21} & \Delta_{22}
\end{array}
\]
and the direct hidden-continuation contribution is
\[
h_q(x)
=
\sum_{(i,j)\in S^2}\Delta_{ij}(q;x).
\]

\begin{theorem}[Local distance reduction]
\label{thm:7.1}
For every \(x\in\Omega'\) and \(q\in S^5\),
\[
\boxed{
h_q(x)=
2\left(
D_C(t,xy,f;b,a,z_R)-D_L(ep_L,ep_R)
\right).
}
\]
Consequently the local direct definition and the Hamming-distance formula
define the same global observables \(rawH\) and \(H_{tot}\).
\end{theorem}

\begin{proof}
By definition, \(D_C(t,xy,f;b,a,z_R)\) is the sum of the nine indicators
\(\delta^C_{ij}(q;x)\), and \(D_L(ep_L,ep_R)\) is the sum of the nine
indicators \(\delta^L_{ij}(q;x)\). Subtracting these two Hamming
distances and multiplying by \(2\) gives exactly the displayed sum of
the nine direct local summands \(\Delta_{ij}\). The symbolic reduction
is recorded by the certificate handles \texttt{S6-RED} and
\texttt{S6-CERT-IF}.
\end{proof}

Then
\[
rawH(x)=\sum_{q\in S^5}\frac{h_q(x)}2,
\qquad
H_{tot}(x)=6\,rawH(x)=3\sum_{q\in S^5}h_q(x),
\]
and
\[
\boxed{
N_-(x)=3\,\#\{q\in S^5:h_q(x)<0\}.
}
\]
Thus purity is the finite system
\[
N_-(x)=0
\quad\Longleftrightarrow\quad
h_q(x)\ge0\quad(q\in S^5).
\]

\subsection{Global frontier theorem}\label{global-frontier-theorem}

The global hidden-continuation theorem is supported by four finite
certificate obligations.

\begin{longtable}[]{@{}
  >{\raggedright\arraybackslash}p{(\columnwidth - 4\tabcolsep) * \real{0.3333}}
  >{\raggedright\arraybackslash}p{(\columnwidth - 4\tabcolsep) * \real{0.3333}}
  >{\raggedright\arraybackslash}p{(\columnwidth - 4\tabcolsep) * \real{0.3333}}@{}}
\toprule\noalign{}
\begin{minipage}[b]{\linewidth}\raggedright
handle
\end{minipage} & \begin{minipage}[b]{\linewidth}\raggedright
obligation
\end{minipage} & \begin{minipage}[b]{\linewidth}\raggedright
result
\end{minipage} \\
\midrule\noalign{}
\endhead
\bottomrule\noalign{}
\endlastfoot
H1 & unrestricted full-landscape range & \(rawH\in[-378,1217]\),
endpoint counts \(8,12\) \\
H2 & pure-frontier exclusion &
\((h_q\ge0\ \forall q)\wedge rawH\ge1171\) is UNSAT \\
H3 & pure-frontier witnesses & PAB and row-complement have
\(rawH=1170\), \(N_-=0\) \\
H4 & signed-cancellation spike & exactly twelve points above the pure
frontier, all with \((rawH,N_-)=(1217,3)\) \\
\end{longtable}

\begin{theorem}[Global frontier of \(\mathcal H\)]
\label{thm:7.2}
On
\(\Omega'=S^{21}\):

\begin{enumerate}
\def\labelenumi{\arabic{enumi}.}
\item
  The unrestricted range is \[
  \boxed{
  \min H_{tot}=-2268,
  \qquad
  \max H_{tot}=7302.
  }
  \] Equivalently, \(rawH\in[-378,1217]\). The endpoint multiplicities
  are \(8\) and \(12\).
\item
  The pure frontier is \[
  \boxed{
  \max\{H_{tot}:N_-=0\}=7020.
  }
  \]
\item
  The region above the pure frontier is exactly the unrestricted maximum
  layer: \[
  \boxed{
  \{H_{tot}>7020\}=\{H_{tot}=7302\}.
  }
  \] This layer has \(12\) points, all with \[
  rawH=1217,
  \qquad H_{tot}=7302,
  \qquad N_-=3,
  \] \[
  H_{pos}=7308,
  \qquad H_{\mathrm{neg\_abs}}=6,
  \qquad h_{\mathrm{loc,min}}=-2,
  \qquad h_{\mathrm{loc,max}}=18.
  \]
\end{enumerate}
\end{theorem}

\begin{proof}
The analytic reduction is Theorem 7.1. H1 covers all
\(3^{21}\) normal-form words and gives the unrestricted range. For the
pure frontier, \(N_-=0\) is the system \(h_q\ge0\) for all \(q\in S^5\),
and \(H_{tot}>7020\) is equivalent to \(rawH\ge1171\). H2 proves \[
(h_q\ge0\ \forall q\in S^5)\wedge rawH\ge1171
\] UNSAT on all \(18540\) depth-9 live nodes: \[
18540/18540\ \mathrm{UNSAT},
\qquad SAT=0,
\qquad UNKNOWN=0.
\] H3 supplies pure witnesses at \(rawH=1170\), and H4 classifies the
full layer above the pure threshold as the twelve impure maximizers.
\end{proof}

\subsection{PAB and row-complement on the pure
frontier}\label{pab-and-row-complement-on-the-pure-frontier}

\begin{corollary}
\label{cor:7.3}
PAB and row-complement satisfy \[
H_{tot}=7020,
\qquad rawH=1170,
\qquad N_-=0.
\] Thus both lie on the global pure hidden-continuation frontier.
\end{corollary}

\begin{proof}
Direct evaluation of \[
\mathrm{PAB}=111111111222222222000,
\qquad
\mathrm{comp}=222222222111111111000
\] gives the displayed values, and Theorem 7.2 proves that no pure point
has larger \(H_{tot}\).
\end{proof}

Part I has now supplied the ambient finite map: normal form, global
associativity, and the hidden-continuation frontier. Part II turns from
atlas data to the intrinsic selector. Its proof uses the access-zero and
diagonal-compression reductions, then the directed-edge absorption
geometry; the frontier theorem supplies the global placement of the two
finalists.

\part{Intrinsic finite selection}

\section{Selection core}\label{selection-core}

The compact uniqueness core is \[
\boxed{
\mathcal H_{acc}=0,
\qquad
H_{diag}=0,
\qquad
\text{diagonal idempotence},
\qquad
\text{pure }C/J.
}
\] The order of the four conditions mirrors the row/column geometry of
\(\Omega'\). Zero hidden-column access makes cross-row output
column-blind. Zero diagonal entropy makes the same-point diagonal sector
constant. Diagonal idempotence fixes that constant to \(d=000\). Pure
\(C/J\) then reads the absorption geometry on the six directed edges
\(M^\times\) and chooses the undrifted continuation/reversal pair.

The longer explanatory chain is \[
\Omega'
\xrightarrow{\mathcal H_{acc}=0,
\ H_{diag}=0}
9\times3
\xrightarrow{a\ne b,
\ d\text{ idempotent}}
6\times1
\xrightarrow{\min\operatorname{Assoc}_{000}}
2\times1
\xrightarrow{\text{pure }C/J}
1.
\] It records the local path-dependence diagnostic that exposes the two
finalists, PAB and row-complement, before the pure \(C/J\) separation.

\section{Diagonal compression}\label{diagonal-compression}

The access-zero theorem has already compressed the cross-row sector to
the nine column-blind rules \(h_{a,b}\). The remaining
information-compression step acts on the same-point diagonal data.

For a diagonal triple \[
d=(d_0,d_1,d_2)\in S^3,
\] let \[
p_s=\frac13\#\{i:d_i=s\}.
\] Define the diagonal Shannon observable by \[
\boxed{
H_{diag}(d)=-\sum_{s\in S}p_s\log_2 p_s.
}
\] The possible levels are:

\begin{longtable}[]{@{}lrr@{}}
\toprule\noalign{}
diagonal type & multiplicities & \(H_{diag}\) \\
\midrule\noalign{}
\endhead
\bottomrule\noalign{}
\endlastfoot
constant & \((3,0,0)\) & \(0\) \\
two-valued & \((2,1,0)\) & \(\log_2 3-\frac23\) \\
permutation & \((1,1,1)\) & \(\log_2 3\) \\
\end{longtable}

\begin{theorem}
\label{thm:9.1}
For diagonal triples, \[
H_{diag}(d)=0
\quad\Longleftrightarrow\quad
 d_0=d_1=d_2.
\] Thus diagonal compression reduces the diagonal sector from \(27\)
choices to \[
000,
\qquad 111,
\qquad 222.
\]
\end{theorem}

\begin{proof}
Entropy of a finite distribution vanishes exactly when
the distribution is concentrated on one symbol. For the empirical
distribution of \((d_0,d_1,d_2)\), this is precisely constancy of the
diagonal triple.
\end{proof}

Together with the access-zero theorem of Section 3, this reduces the
selection domain to \(9\times3\).

\section{Anchors and local
path-dependence}\label{anchors-and-local-path-dependence}

\subsection{Nondegenerate anchors}\label{nondegenerate-anchors}

\begin{proposition}
\label{prop:10.1}
Among column-blind cross-row rules
\(h_{a,b}\), the condition \[
a\ne b
\] leaves exactly six rules.
\end{proposition}

\begin{proof}
There are nine pairs \((a,b)\in S^2\). The three
diagonal pairs \((a,a)\) do not distinguish the two cross-row
representatives. Removing them leaves \(3^2-3=6\) nondegenerate rules.
\end{proof}

\begin{proposition}
\label{prop:10.2}
After diagonal compression, diagonal
idempotence selects the unique constant diagonal \[
d=000.
\]
\end{proposition}

\begin{proof}
The compressed diagonal sector is \(\{000,111,222\}\).
Diagonal idempotence requires the product \((r,r)\cdot(r,r)\) to return
\((r,r)\), which fixes the constant diagonal \(000\) in the normalized
coordinates.
\end{proof}

\subsection{\texorpdfstring{Local associativity at
\(d=000\)}{Local associativity at d=000}}\label{local-associativity-at-d000}

\begin{proposition}
\label{prop:10.3}
With \(d=000\) fixed, the six nontrivial
column-blind rules have local associativity values \[
\operatorname{Assoc}_{000}=219
\] for exactly two rules, and \[
\operatorname{Assoc}_{000}=273
\] for the remaining four. The two local minima are PAB and
row-complement.

The local-minimum profile is \[
(19,0,0,54,0),
\] whereas the four discarded rules have profile \[
(19,9,9,54,0).
\] The difference comes from the additional \(RRS\) and \(RSR\)
coincidences in the discarded rules.
\end{proposition}

\begin{proof}
Proposition 4.3 gives \[
\operatorname{Assoc}_{000}(h_{a,b})
=
3\left(73+18[a=0]+18[b=0]+54[a=b]\right).
\] On the nondegenerate stratum \(a\ne b\), this value is minimized
exactly when neither anchor is \(0\), i.e.~at \((a,b)=(1,2)\) and
\((2,1)\). These are PAB and row-complement, and both have value
\(219\). The four remaining nondegenerate pairs contain one zero anchor
and have value \(273\). The certificate table
\texttt{l2\_path\_dependence\_table.csv} and verifier
\texttt{verify\_l2\_selection.py} record the same six-row audit.
\end{proof}

This step is local and conditional. It is applied only after information
compression, nondegenerate anchoring, and diagonal anchoring.

\section{Directed-edge geometry}\label{directed-edge-geometry}

Let \[
M^\times=\{(r,c)\in M:r\ne c\}
\] be the six-point directed-edge set. Each \((r,c)\in M^\times\) is
interpreted as the directed edge from \(r\) to \(c\) in the triangle on
vertex set \(S\).

Define \[
C(r,c)=(c,\overline{rc}),
\qquad
J(r,c)=(c,r).
\] The map \(C\) advances an oriented edge through the third vertex, and
\(J\) reverses orientation.

\begin{theorem}
\label{thm:11.1}
On \(M^\times\), \[
C^3=1,
\qquad
J^2=1,
\qquad
JCJ=C^{-1}.
\] Thus \(C\) and \(J\) generate the dihedral geometry of the directed
triangle.
\end{theorem}

\begin{proof}
Applying \(C\) three times traverses \[
(r,c)\mapsto(c,\overline{rc})\mapsto(\overline{rc},r)\mapsto(r,c).
\] Applying \(J\) twice reverses orientation twice. Finally, \[
JCJ(r,c)=JC(c,r)=J(r,\overline{rc})=(\overline{rc},r)=C^{-1}(r,c).
\] The table \texttt{l3\_CJ\_finite\_relations.csv} audits the finite
relation checks.
\end{proof}

\section{\texorpdfstring{Pure \(C/J\) and the selection
theorem}{Pure C/J and the selection theorem}}\label{pure-cj-and-the-selection-theorem}

\subsection{Absorption and absorption
transitions}\label{absorption-and-absorption-transitions}

For \(x\in M^\times\), define the absorption set \[
\operatorname{Abs}(x)=\{y\in M^\times\setminus\{x\}:x\cdot y=x\}.
\] For a column-blind rule \(h_{a,b}\) at \(d=000\), same-row right
arguments never absorb an off-diagonal \(x\). Thus absorption is
governed by the cross-row equation \[
h(r_x,r_y)=c_x.
\]

A candidate has a two-transition absorption geometry if every
\(x\in M^\times\) has exactly two absorbed neighbours. In that case the
two absorbed-neighbour maps define \(\operatorname{AbsTrans}\). The pure
\(C/J\) criterion requires \[
\operatorname{AbsTrans}=\{C,J\}
\] as literal maps on \(M^\times\), with no \(C^{\pm1}\)-drift.

\subsection{The final two candidates}\label{the-final-two-candidates}

After local path-dependence, the two remaining rules are PAB and
row-complement.

\begin{lemma}
\label{lem:12.1}
For PAB, \[
\operatorname{AbsTrans}_{PAB}=\{C,J\}.
\] For row-complement, \[
\operatorname{AbsTrans}_{comp}=\{C^{-1},C^{-1}J\}.
\]
\end{lemma}

\begin{proof}
Let \(x=(r,c)\in M^\times\).

For PAB, \(h(r_x,r_y)=r_y\). The absorption equation is \[
r_y=c.
\] Hence the two absorbed neighbours of \((r,c)\) are the two
off-diagonal points in row \(c\): \[
(c,r),
\qquad
(c,\overline{rc}).
\] These are precisely \[
J(r,c)=(c,r),
\qquad
C(r,c)=(c,\overline{rc}).
\] So \(\operatorname{AbsTrans}_{PAB}=\{C,J\}\).

For row-complement, \(h(r_x,r_y)=\overline{r_xr_y}\). The absorption
equation is \[
\overline{r r_y}=c,
\] so \[
r_y=\overline{rc}.
\] The two absorbed neighbours are therefore \[
(\overline{rc},r),
\qquad
(\overline{rc},c),
\] which are \[
C^{-1}(r,c)=(\overline{rc},r),
\qquad
C^{-1}J(r,c)=(\overline{rc},c).
\] Thus \(\operatorname{AbsTrans}_{comp}=\{C^{-1},C^{-1}J\}\).
\end{proof}

\subsection{Final selector among the two local
finalists}\label{final-selector-among-the-two-local-finalists}

\begin{theorem}
\label{thm:12.2}
Among the final two candidates, the pure \(C/J\)
directed-edge criterion selects PAB.
\end{theorem}

\begin{proof}
By Lemma 12.1, PAB realizes the undrifted pair
\(\{C,J\}\). Row-complement realizes \(\{C^{-1},C^{-1}J\}\), whose
continuation and reversal are drifted by \(C^{-1}\). The pure criterion
accepts only the literal undrifted pair. The audit table
\texttt{l2\_pure\_CJ\_survivors.csv} records the same comparison.
\end{proof}

\subsection{Compact uniqueness core}\label{compact-uniqueness-core}

\begin{theorem}
\label{thm:12.3}
The four conditions \[
\mathcal H_{acc}=0,
\qquad
H_{diag}=0,
\qquad
\text{diagonal idempotence},
\qquad
\text{pure }C/J
\] select PAB uniquely inside \(\Omega'\).
\end{theorem}

\begin{proof}
The first condition gives column-blindness by Theorem
3.1. Hence the cross-row rule has the form \[
h(r,r+1)=r+a,
\qquad
h(r,r+2)=r+b
\] for a unique pair \((a,b)\in S^2\). The second condition gives a
constant diagonal by Theorem 9.1, and diagonal idempotence fixes this
constant to \(d=000\) by Proposition 10.2.

It remains to impose pure \(C/J\). Let \(x=(r,c)\in M^\times\) and write
\(\delta=c-r\in\{1,2\}\). Same-row right arguments do not absorb \(x\)
at \(d=000\), so absorption is governed by the cross-row equation \[
x\cdot y=x
\quad\Longleftrightarrow\quad
h(r,r_y)=c.
\] For the two possible cross rows this becomes \[
r_y=r+1\Rightarrow a=\delta,
\qquad
r_y=r+2\Rightarrow b=\delta.
\]

The pure \(C/J\) criterion requires the two absorbed neighbours of
\(x=(r,c)\) to be exactly \[
J(r,c)=(c,r),
\qquad
C(r,c)=(c,\overline{rc}).
\] Both lie in row \(c\). Therefore, when \(\delta=1\), absorption must
occur in row \(r+1=c\), forcing \(a=1\). When \(\delta=2\), absorption
must occur in row \(r+2=c\), forcing \(b=2\). Thus \[
(a,b)=(1,2).
\] This is precisely the PAB cross-row rule, and \(d=000\) has already
been fixed. Hence the four conditions select PAB uniquely.
\end{proof}

The audit table \texttt{l2\_minimal\_selector\_sets.csv} records the
same compact core as a finite audit of the survivor condition.

\subsection{Full explanatory selection
chain}\label{full-explanatory-selection-chain}

\begin{proposition}
\label{prop:12.4}
The finite explanatory chain
\[
\Omega'
\xrightarrow{\mathcal H_{acc}=0,
\ H_{diag}=0}
9\times3
\xrightarrow{a\ne b,
\ d\text{ idempotent}}
6\times1
\xrightarrow{\min\operatorname{Assoc}_{000}}
2\times1
\xrightarrow{\text{pure }C/J}
1
\]
has unique survivor \(\mathrm{PAB}\).
\end{proposition}

\begin{proof}
The first arrow is Theorems 3.1 and 9.1. The second
arrow is Propositions 10.1 and 10.2. The third arrow is Proposition
10.3. The final arrow is Theorem 12.2. The audit table
\texttt{l2\_selection\_chain.csv} records the survivor count at each
stage, and the compact L2 verifier checks the displayed chain.
\end{proof}

This chain is logically subsumed by Theorem~\ref{thm:12.3}. After
column-blindness, diagonal compression, and diagonal idempotence, the
pure \(C/J\) condition alone forces \((a,b)=(1,2)\). The
\(\min\operatorname{Assoc}_{000}\) arrow is retained only as a didactic
intermediate: it exposes the PAB/row-complement ambiguity before the
directed-edge orientation criterion separates them.

\part{Linearization after selection}

\section{The selected PAB operation}\label{the-selected-pab-operation}

\subsection{PAB multiplication and ground
fields}\label{pab-multiplication-and-ground-fields}

The selected multiplication is \[
(r_1,c_1)\cdot(r_2,c_2)=
\begin{cases}
(r_1,r_2), & r_1\ne r_2,\\[2mm]
(r_1,\overline{c_1c_2}), & r_1=r_2,\ c_1\ne c_2,\\[2mm]
(r_1,r_1), & r_1=r_2,\ c_1=c_2.
\end{cases}
\] The three cases are cross-row multiplication, same-row off-diagonal
Steiner multiplication, and diagonal anchoring.

The finite magma statements are set-theoretic. When the carrier is
linearized, let \(k\) be a ground field and let \(k[M]\) be the vector
space with basis \(e_x\), \(x\in M\). Rank and image statements for
individual finite maps are valid over any ground field. The finite
composition monoid and term-function cardinalities are
field-independent; linear span dimensions are dimensions over the chosen
\(k\). The left-translation span in Theorem 14.2 is
characteristic-independent by the displayed normal-form proof. Spectral,
Jordan, and fiber-algebra statements involving the factors \(t-1\),
\(t+1\), or the idempotents with coefficients \(1/2\) are stated under
\[
\operatorname{char} k\ne2.
\] The compact Lie and symplectic bridge statements are over
\(\mathbb R\) unless explicitly stated otherwise.

The carrier vector space admits the matrix-unit identification \[
k[M]\cong\operatorname{Mat}_3(k),
\qquad
 e_{(r,c)}\leftrightarrow E_{rc}.
\] This identification is used for linear and bridge structures. PAB
multiplication remains the finite multiplication displayed above.

\subsection{Row-owner principle}\label{row-owner-principle}

For a nonempty term \(t(x_1,\ldots,x_n)\), define its \textbf{owner}
\(\operatorname{ow}(t)\) to be the leftmost variable appearing in the
term tree.

\begin{lemma}[Row-owner lemma]
\label{lem:13.1}
For every PAB term
\(t(x_1,\ldots,x_n)\), \[
row\bigl(t(x_1,\ldots,x_n)\bigr)
=
row\bigl(\operatorname{ow}(t)\bigr).
\]
\end{lemma}

\begin{proof}
The assertion is immediate for a variable. If
\(t=u\cdot v\), then the left-row output axiom gives \[
row(t)=row(u\cdot v)=row(u).
\] By induction, \(row(u)\) is the row of the leftmost variable of
\(u\), which is also the leftmost variable of \(u\cdot v\).
\end{proof}

This lemma is the structural source of the owner coordinate in the
binary term algebra of Section 18 and of the row-target law for
compositions of left translations.

\section{Left regular operators}\label{left-regular-operators}

For \(x\in M\), define the left translation \[
L_x(y)=x\cdot y.
\]

\subsection{Row-fiber normal form}\label{row-fiber-normal-form}

\begin{proposition}
\label{prop:14.1}
For \(x=(r,c)\), the left translation has the
explicit form \[
L_{(r,c)}(s,d)=
\begin{cases}
(r,s), & s\ne r,\\[1mm]
(r,r), & s=r,\ d=c,\\[1mm]
(r,\overline{cd}), & s=r,\ d\ne c.
\end{cases}
\] In particular, \[
L_{(r,c)}(M)=F_r.
\]
\end{proposition}

\begin{proof}
If \(s\ne r\), the PAB cross-row rule gives
\((r,c)\cdot(s,d)=(r,s)\). If \(s=r\) and \(d=c\), diagonal anchoring
gives \((r,r)\). If \(s=r\) and \(d\ne c\), the same-row Steiner rule
gives \((r,\overline{cd})\). These three cases also show that the image
contains all three points of \(F_r\).
\end{proof}

\subsection{Rank and span}\label{rank-and-span}

\begin{theorem}
\label{thm:14.2}
For every \(x\in M\), \[
\operatorname{rank}L_x=3,
\qquad
\operatorname{Im}L_x=F_{row(x)}.
\] Moreover, \[
\dim\operatorname{span}\{L_x:x\in M\}=9.
\]
\end{theorem}

\begin{proof}
Proposition 14.1 shows that each \(L_x\) maps onto the
three-point row fiber \(F_{row(x)}\). After linearization, the three
basis vectors of that row fiber occur among the columns of the matrix of
\(L_x\), so the rank is \(3\).

It remains to prove the span dimension. Operators with different target
rows have disjoint codomain support, so their spans are independent. Fix
a row \(r\). The three operators \(L_{(r,c)}\), \(c\in S\), are
independent because their restrictions to \(kF_r\) are, after a
relabelling of \(S\), the three matrices \[
\begin{pmatrix}1&0&0\\0&0&1\\0&1&0\end{pmatrix},
\qquad
\begin{pmatrix}0&1&1\\0&0&0\\1&0&0\end{pmatrix},
\qquad
\begin{pmatrix}0&1&1\\1&0&0\\0&0&0\end{pmatrix},
\] and these are visibly linearly independent. Thus each row contributes
dimension \(3\), and the total span dimension is \(3\cdot3=9\). The
table \texttt{l3\_operator\_spectrum\_rank.csv} audits the nine-row
calculation.
\end{proof}

\subsection{Spectral dichotomy}\label{spectral-dichotomy}

Assume \(\operatorname{char}k\ne2\).

\begin{theorem}
\label{thm:14.3}
If \(x=(r,r)\) is diagonal, then \[
\mu_{L_x}(t)=t(t-1)(t+1),
\] and \(L_x\) is semisimple. If \(x=(r,c)\) with \(r\ne c\), then \[
\mu_{L_x}(t)=t^2(t-1)(t+1),
\] with zero-primary Jordan structure \[
J_2(0)\oplus J_1(0)^{\oplus5}.
\]
\end{theorem}

\begin{proof}
Write \(V=k[M]\) and \(V_r=kF_r\). Every \(L_{(r,c)}\)
maps \(V\) onto \(V_r\). Hence all nonzero spectral data are controlled
by the restriction \[
T_{r,c}=L_{(r,c)}|_{V_r}:V_r\to V_r.
\]

If \(c=r\), the map \(T_{r,r}\) fixes \(e_{(r,r)}\) and swaps the two
off-diagonal basis vectors in \(V_r\). Thus its minimal polynomial is
\((t-1)(t+1)\). Since \(T_{r,r}\) is invertible,
\(V=\ker L_{(r,r)}\oplus V_r\), and \(L_{(r,r)}\) has minimal polynomial
\(t(t-1)(t+1)\) and no nilpotent zero-primary extension.

If \(c\ne r\), the restriction \(T_{r,c}\) has eigenvalues \(0,1,-1\) on
\(V_r\). Its kernel is one-dimensional, so \[
\operatorname{Im}L_{(r,c)}\cap\ker L_{(r,c)}
=
\ker T_{r,c}
\] is one-dimensional. This produces one length-two Jordan block at
eigenvalue \(0\), while the nonzero \(\pm1\) part remains semisimple.
Therefore the minimal polynomial is \(t^2(t-1)(t+1)\). The table
\texttt{l3\_operator\_spectrum\_rank.csv} audits the spectral/Jordan
calculation for all nine left translations.
\end{proof}

\section{Composition envelope and directed-edge
representation}\label{composition-envelope-and-directed-edge-representation}

\subsection{Fixed-row normal forms for
compositions}\label{fixed-row-normal-forms-for-compositions}

Let \(\operatorname{Mon}^+\langle L_x\rangle\) be the positive
composition monoid generated by the nine left translations. For
\(r\in S\), let \(\mathcal W_r\) be the component of positive
compositions whose output lies in \(F_r\).

For fixed \(r\), define \[
\tau_{r,c}:S\to S,
\qquad
\tau_{r,c}(d)=
\begin{cases}
r, & d=c,\\
\overline{cd}, & d\ne c.
\end{cases}
\] This is the restriction of \(L_{(r,c)}\) to \(F_r\), after
identifying \(F_r\) with its column set. Let \[
\mathcal P_r=\langle \tau_{r,c}:c\in S\rangle.
\] In normalized row \(0\), \[
\mathcal P_0=
\{012,021,200,100,011,022\}.
\]

For \(\psi\in\mathcal P_r\) and \(c\in S\), set \[
N_{r,\psi,c}(s,d)=
\begin{cases}
(r,\psi(\tau_{r,c}(d))), & s=r,\\
(r,\psi(s)), & s\ne r.
\end{cases}
\] Also let \(K_{r,t}\) be the constant map to \((r,t)\).

\begin{proposition}[Fixed-row composition normal form]
\label{prop:15.1}
For every
\(r\in S\), \[
\boxed{
\mathcal W_r
=
\{N_{r,\psi,c}:\psi\in\mathcal P_r,
\ c\in S\}
\cup
\{K_{r,t}:t\in S\}.
}
\] The first set contains \(14\) distinct nonconstant maps, the second
contains \(3\) constants, and therefore \[
|\mathcal W_r|=17.
\]
\end{proposition}

\begin{proof}
A positive composition has a leftmost generator, and by
the row-owner lemma its target row is the row \(r\) of that generator.
If every generator in the word has target row \(r\), the rightmost
generator is some \(L_{(r,c)}\); it applies \(\tau_{r,c}\) on \(F_r\)
and sends every other row \(F_s\) to column \(s\). The remaining left
prefix is an element \(\psi\in\mathcal P_r\), yielding \(N_{r,\psi,c}\).
Conversely each such form is realized by a word for \(\psi\) followed by
\(L_{(r,c)}\).

If a generator to the right has target row \(s\ne r\), the leftmost
row-\(r\) generator sees a cross-row input and the word becomes
constant; the three constants are realized by two-letter words and
row-fiber prefixes. In row \(0\), the six-by-three list
\(\mathcal P_0\times S\) has four twofold collisions and no other
collisions, leaving \(18-4=14\) nonconstant forms. Translation gives the
same count for all rows.
\end{proof}

\subsection{Composition envelope
theorem}\label{composition-envelope-theorem}

\begin{theorem}
\label{thm:15.2}
The positive monoid generated by the left
translations has size \[
|\operatorname{Mon}^+\langle L_x\rangle|=51.
\] More precisely, \[
\operatorname{Mon}^+\langle L_x\rangle
=
\mathcal W_0\sqcup\mathcal W_1\sqcup\mathcal W_2,
\qquad
|\mathcal W_r|=17.
\] Its linear envelope decomposes as \[
\mathcal E=W_0\oplus W_1\oplus W_2,
\qquad
\dim W_r=9,
\qquad
\dim\mathcal E=27,
\] and satisfies \[
W_rW_s\subseteq W_r.
\]
\end{theorem}

\begin{proof}
Proposition 15.1 gives \(|\mathcal W_r|=17\), and the
three target-row components are disjoint. Hence the positive composition
monoid has \(3\cdot17=51\) elements. If \(A\in\mathcal W_r\) and
\(B\in\mathcal W_s\), then \(A\circ B\) has target row \(r\), proving
\(W_rW_s\subseteq W_r\).

For the span dimension, it suffices by symmetry to work in row \(0\).
The certificate exposes nine normalized maps whose one-hot
linearizations have a determinant \(-1\) minor, and reduces the
remaining eight maps to this basis by integer relations. Thus
\(\dim W_0=9\), and translating rows gives \(\dim W_r=9\) for all \(r\).
Disjoint codomain supports give the direct sum and
\(\dim\mathcal E=27\). The tables \texttt{l3\_composition\_envelope.csv}
and \texttt{l3\_composition\_product\_law.csv} audit the normal forms,
basis reductions, component counts, and row-target product law.
\end{proof}

\subsection{The five-dimensional row-fiber
algebra}\label{the-five-dimensional-row-fiber-algebra}

Assume \(\operatorname{char}k\ne2\). For each row \(r\), define \[
A_{5,r}=\langle L_{(r,c)}|_{kF_r}:c\in S\rangle\subseteq\operatorname{End}(kF_r).
\]

\begin{proposition}
\label{prop:15.3}
The algebra \(A_{5,r}\) has dimension \(5\).
For \(r=0\), with \[
B_1=L_{(0,0)}|_{kF_0},
\qquad
B_2=L_{(0,1)}|_{kF_0},
\qquad
B_3=L_{(0,2)}|_{kF_0},
\qquad
B_5=B_2^2,
\] one has the basis \[
B_1,
\quad B_2,
\quad B_3,
\quad I,
\quad B_5.
\] Furthermore, \[
Z(A_{5,r})=kI,
\] \(A_{5,r}\) has a two-dimensional square-zero radical, and \[
A_{5,r}/\operatorname{Rad}(A_{5,r})\cong k^3.
\]
\end{proposition}

\begin{proof}
By \(\sigma\)-equivariance it is enough to check
\(r=0\). The multiplication table of the five displayed basis elements
closes and has dimension \(5\). The retained fiber-algebra certificate
identifies the center as \(kI\), the radical as a two-dimensional
square-zero ideal, and the quotient as \(k^3\). The supporting tables
are \texttt{l3\_fiber\_algebra\_5dim\_multiplication.csv} and
\texttt{l3\_fiber\_algebra\_5dim\_structure.csv}, checked by
\texttt{verify\_l3\_fiber\_algebra\_5dim.py}.
\end{proof}

\subsection{Directed-edge
representation}\label{directed-edge-representation}

\begin{theorem}
\label{thm:15.4}
On the six directed edges, \[
\mathbb R[M^\times]\cong\mathbf1\oplus\mathrm{sign}\oplus2\,\mathrm{std}.
\] For PAB the absorption matrix is \[
A_{abs}=(P_C+P_J)^\top
\] and has spectrum \[
\{2,0,0,0,-1,-1\}.
\]
\end{theorem}

\begin{proof}
The six directed edges carry the permutation
representation of the directed triangle. Its character decomposition is
\(\mathbf1\oplus\mathrm{sign}\oplus2\mathrm{std}\). PAB absorption is
generated by the two pure transitions \(C\) and \(J\), so the absorption
matrix is \((P_C+P_J)^\top\). The tables
\texttt{l3\_absorption\_transitions.csv} and
\texttt{l3\_absorption\_spectrum.csv} audit the transition action and
spectrum.
\end{proof}

\section{Endomorphisms, intertwiners, and row
contraction}\label{endomorphisms-intertwiners-and-row-contraction}

\subsection{Endomorphisms and
intertwiners}\label{endomorphisms-and-intertwiners}

For a map \(\rho:S\to S\), write \[
\varphi_\rho(r,c)=(\rho(r),\rho(c))
\] when this formula defines a magma endomorphism. For \(t\in S\), write
\[
e_t(r,c)=(t,t)
\] for the collapse to the diagonal point \((t,t)\).

For a row \(r\), write the induced column operation on \(S\) as \[
c\star_r d=
\begin{cases}
r, & c=d,\\
\overline{cd}, & c\ne d.
\end{cases}
\]

\begin{lemma}[row-fiber homomorphisms]
\label{lem:16.1}
Let \(\phi:S\to S\)
satisfy \(\phi(r)=t\) and \[
\phi(c\star_r d)=\phi(c)\star_t\phi(d)
\qquad(c,d\in S).
\] Then either \(\phi\) is the constant map to \(t\), or \(\phi\) is a
permutation of \(S\) sending \(r\) to \(t\).
\end{lemma}

\begin{proof}
By translating symbols, it is enough to take \(r=t=0\).
Write \[
\phi(0)=0,
\qquad
\phi(1)=u,
\qquad
\phi(2)=v.
\] The equations \(0\star_0 1=2\) and \(0\star_0 2=1\) give \[
v=0\star_0 u,
\qquad
u=0\star_0 v.
\] These two equations have exactly the three solutions \[
(u,v)=(0,0),
\qquad
(u,v)=(1,2),
\qquad
(u,v)=(2,1).
\] They are respectively the constant map and the two permutations
fixing \(0\), and each directly preserves \(\star_0\). Translating back
gives the claim.
\end{proof}

\begin{theorem}
\label{thm:16.2}
The endomorphism monoid of the selected PAB magma
is \[
\operatorname{End}(M,\cdot)
=
\{\varphi_\rho:\rho\in\operatorname{Sym}(S)\}\cup\{e_0,e_1,e_2\}.
\]
\end{theorem}

\begin{proof}
The idempotents of PAB are exactly the diagonal points
\((r,r)\), since \[
(r,c)^2=(r,r)
\] and this equals \((r,c)\) only for \(c=r\). Hence an endomorphism
\(f\) sends diagonal points to diagonal points. Write \[
f(r,r)=(\rho(r),\rho(r)).
\] For \(x=(r,c)\in F_r\), the identity \(x^2=(r,r)\) gives \[
f(x)^2=f(r,r)=(\rho(r),\rho(r)),
\] so \(f(x)\in F_{\rho(r)}\). Therefore \[
f(r,c)=(\rho(r),\phi_r(c))
\] for maps \(\phi_r:S\to S\) satisfying \(\phi_r(r)=\rho(r)\). Applying
the endomorphism equation to pairs inside \(F_r\) shows that \(\phi_r\)
is a row-fiber homomorphism. By Lemma 16.1, each \(\phi_r\) is either
the constant map to \(\rho(r)\) or a permutation of \(S\) sending \(r\)
to \(\rho(r)\).

Now impose cross-row multiplication. If \(r\ne s\), then \[
f((r,c)\cdot(s,d))=f(r,s)=(\rho(r),\phi_r(s)).
\] If \(\rho(r)\ne\rho(s)\), the right side of the endomorphism equation
is cross-row, so \[
f(r,c)\cdot f(s,d)=(\rho(r),\rho(s)),
\] and hence \[
\phi_r(s)=\rho(s).
\] If \(\rho\) is injective, this holds for every \(r\ne s\), and
together with \(\phi_r(r)=\rho(r)\) gives \(\phi_r=\rho\) for all \(r\).
Thus \(\rho\in\operatorname{Sym}(S)\) and \(f=\varphi_\rho\).

It remains to exclude nonconstant maps with noninjective \(\rho\).
Choose \(r\ne s\) with \(\rho(r)=\rho(s)=t\). The cross-row equation
becomes \[
\phi_r(s)=\phi_r(c)\star_t\phi_s(d)
\qquad(c,d\in S).
\] For fixed \(d\), the map \(a\mapsto a\star_t\phi_s(d)\) is not
constant on \(S\). Therefore \(\phi_r\) cannot be a permutation, so
Lemma 16.1 forces \(\phi_r\) to be constant to \(t\). Symmetrically,
\(\phi_s\) is constant to \(t\).

Let \(u\) be the third row. If \(\rho(u)\ne t\), comparing row \(u\)
with the collapsed rows \(r\) and \(s\) gives \[
\phi_u(r)=t,
\qquad
\phi_u(s)=t.
\] This contradicts Lemma 16.1: \(\phi_u\) is neither a permutation nor
the constant map to \(\rho(u)\). Hence \(\rho(u)=t\), and then applying
the preceding collapsed-pair argument to \(u\) and \(r\) forces
\(\phi_u\) to be constant to \(t\). Thus \(f=e_t\).

Conversely, every \(\varphi_\rho\) with \(\rho\in\operatorname{Sym}(S)\)
preserves the three cases in the PAB multiplication, and every collapse
\(e_t\) is an endomorphism because
\(e_t(x\cdot y)=(t,t)=e_t(x)\cdot e_t(y)\). This gives exactly six
permutation endomorphisms and three collapses.
\end{proof}

For an endomorphism \(\varphi\), let \(\operatorname{Hom}_\varphi\)
denote the linear space of intertwiners satisfying \[
T L_x=L_{\varphi(x)}T
\qquad(x\in M).
\] For automorphisms, \[
\operatorname{Hom}_{\varphi_\rho}=kP_\rho,
\qquad
\dim=1.
\] For collapses, \[
\operatorname{Hom}_{e_r}=\{u\varepsilon:u\in\operatorname{span}\{d_r,o_r\}\},
\qquad
\dim=2.
\] Therefore \[
\sum_\varphi\dim\operatorname{Hom}_\varphi=6\cdot1+3\cdot2=12.
\] The displayed intertwiner formulas are checked by substitution into
the intertwining equations; the upper bounds and dimensions are audited
in \texttt{l3\_intertwiner\_audit.csv}. Thus the table is an audit of
the linear equations after the endomorphisms have been classified
structurally.

\subsection{Cross-row row contraction}\label{cross-row-row-contraction}

\begin{proposition}
\label{prop:16.3}
Under the carrier identification
\(k[M]\cong\operatorname{Mat}_3(k)\), the row-contraction extension of
the cross-row PAB rule is \[
\mu_{cross}(A,B)=AJ_3B^\top=(A\mathbf1)(B\mathbf1)^\top.
\] On matrix units with cross rows, \[
E_{ij}J_3E_{kl}^\top=E_{ik}.
\]
\end{proposition}

\begin{proof}
Since \(J_3=\mathbf1\mathbf1^\top\), \[
E_{ij}J_3E_{kl}^\top
=(e_ie_j^\top)(\mathbf1\mathbf1^\top)(e_le_k^\top)
=e_i(e_j^\top\mathbf1)(\mathbf1^\top e_l)e_k^\top
=E_{ik}.
\] For \(i\ne k\), this is exactly the cross-row PAB rule
\((i,j)\cdot(k,l)=(i,k)\). The table
\texttt{l3\_mat3\_bridge\_table.csv} audits the basis-level check and
the coincidences with ordinary matrix multiplication.
\end{proof}

With the row-sum augmentation \(\varepsilon(A)=A\mathbf1\), \[
\mu_{cross}(A,B)=\varepsilon(A)\varepsilon(B)^\top.
\] Thus the matrix bridge expresses the same column-blindness selected
by \(\mathcal H_{acc}=0\): cross-row multiplication factors through row
sums. If one extracts only the cross-row branch as a summand of the full
bilinear PAB product on arbitrary matrices, the corresponding expression
is \[
\pi_{off}\!\left(AJ_3B^\top\right).
\] The same-row PAB branch is the Steiner/diagonal rule on row fibers.

\section{Hidden-continuation operator
identity}\label{hidden-continuation-operator-identity}

The local hidden-continuation identity of Section 7 has an operator
interpretation in the left-regular envelope. For a normalized triple
indexed by \(q\), the continuation distance \(D_C\) compares the two
stepwise computations, while \(D_L\) compares the two collapsed left
translations.

\begin{theorem}
\label{thm:17.1}
For every local index \(q\in S^5\), \[
\frac{h_q}{2}=D_C(t,xy,f;b,a,z_R)-D_L(ep_L,ep_R).
\] Equivalently, in the left-regular language this is the finite
Hamming-distance contrast between \[
L_{xy}L_z
\quad\text{and}\quad
L_xL_{yz}
\] on the incremental side, and \[
L_{(xy)z}
\quad\text{and}\quad
L_{x(yz)}
\] on the collapsed side. All four operators lie in the same row-target
component \(W_{row(x)}\).

For PAB, \[
I_{tot}=8904,
\qquad
B_{tot}=1884,
\qquad
H_{tot}=7020,
\qquad
N_-=0.
\] The row-complement has the same pure hidden-continuation profile.
\end{theorem}

\begin{proof}
The first statement is Theorem 7.1 rewritten in operator
language. Row-target compatibility follows from \(W_rW_s\subseteq W_r\).
The bridge audit \texttt{l3\_H\_bridge\_pab\_comp.csv} records the PAB
and row-complement comparison; the finite check verifies that both have
\(H_{tot}=7020\) and \(N_-=0\). The global frontier status follows from
Theorem 7.2.
\end{proof}

The structural reading is one-sided path information. At the pure
frontier, every local contribution is nonnegative: stepwise continuation
never carries less finite distinction than the collapsed
left-translation comparison. The only way to exceed the pure frontier is
the twelve-point impure spike of Theorem 7.2, where signed local
cancellation is present.

\section{Binary term algebra}\label{binary-term-algebra}

\subsection{Binary sectors and normal
signatures}\label{binary-sectors-and-normal-signatures}

A binary term function of PAB is a function \(M^2\to M\) realized by a
term in variables \(x,y\). The set of all such functions modulo equality
is denoted \(T_2(\mathrm{PAB})\).

By the row-owner lemma, every binary term has an owner variable, the
leftmost variable in its term tree: \[
\operatorname{ow}(f)\in\{x,y\}.
\] Order the two inputs as \[
(o,n)=(\operatorname{owner},\operatorname{nonowner}).
\] On the same-row sector, normalize the common row to \(0\) and define
\[
s_f(i,j)=\operatorname{col}\bigl(f(o=(0,i),n=(0,j))\bigr),
\qquad i,j\in S.
\] On the cross-row sector, normalize the owner row to \(0\) and the
nonowner row to \(1\). The output column is independent of the nonowner
column; define \[
u_f(i)=\operatorname{col}\bigl(f(o=(0,i),n=(1,j))\bigr),
\qquad i\in S.
\] The opposite cross-row direction is determined by \[
u_f^{-}(i)=-u_f(-i).
\]

\begin{proposition}[Binary sector reduction]
\label{prop:18.1}
Every binary term
function has a normal signature \[
\Theta(f)=\bigl(\operatorname{ow}(f),s_f,u_f\bigr).
\] The same-row signature \(s_f:S^2\to S\) and cross unary signature
\(u_f:S\to S\) determine the term function on all row sectors.
\end{proposition}

\begin{proof}
The owner coordinate is Lemma 13.1. Variables have the
displayed sector form. If \(f=g\cdot h\), the owner of \(f\) is the
owner of \(g\). On same-row inputs, both subterms remain in input rows,
so the output column is a function of the two normalized columns. On
cross-row inputs, the right subterm is either in the nonowner row,
giving a cross-row output independent of the nonowner column, or in the
owner row, where the same-row PAB rule combines two unary owner-column
signatures. The opposite cross-row direction follows by the
PAB-compatible relabelling fixing the owner row and exchanging the two
nonowner directions.
\end{proof}

\subsection{Exposed finite factors}\label{exposed-finite-factors}

The realized same-row signatures form a \(21\)-element set
\(\mathcal S_{21}\). The realized cross unary signatures form a
\(15\)-element set \(\mathcal U_{15}\). The factorization index lists
every element of both sets and gives realizing terms for the triples
below. Thus the candidate binary normal forms are \[
\{x,y\}\times\mathcal S_{21}\times\mathcal U_{15}.
\] The finite transition tables record the induced pointwise PAB
multiplication on the \(21\)-state and \(15\)-state factors.

\subsection{Binary realization
certificate}\label{binary-realization-certificate}

\begin{lemma}[Binary term certificate protocol]
\label{lem:18.2}
The binary term
theorem is supported by four finite obligations.

\begin{longtable}[]{@{}
  >{\raggedright\arraybackslash}p{(\columnwidth - 2\tabcolsep) * \real{0.5000}}
  >{\raggedright\arraybackslash}p{(\columnwidth - 2\tabcolsep) * \real{0.5000}}@{}}
\toprule\noalign{}
\begin{minipage}[b]{\linewidth}\raggedright
obligation
\end{minipage} & \begin{minipage}[b]{\linewidth}\raggedright
content
\end{minipage} \\
\midrule\noalign{}
\endhead
\bottomrule\noalign{}
\endlastfoot
T2-SIG & the factorization index exposes exactly \(\mathcal S_{21}\) and
\(\mathcal U_{15}\) \\
T2-REAL & every
\((o,s,u)\in\{x,y\}\times\mathcal S_{21}\times\mathcal U_{15}\) has a
realizing term \\
T2-TRANS & pointwise multiplication closes on the two factor transition
systems \\
T2-MULT & the full \(630\times630\) multiplication table is
reconstructed from owner and factor transitions \\
\end{longtable}

The analytic part is the sector reduction and injectivity of \(\Theta\);
the certificate supplies realization and closure of the finite factors.
\end{lemma}

\subsection{Binary normal-form
theorem}\label{binary-normal-form-theorem}

\begin{theorem}
\label{thm:18.3}
The binary term-function algebra of PAB is finite
and has size \[
|T_2(\mathrm{PAB})|=630.
\] More precisely, \[
\boxed{
T_2(\mathrm{PAB})
\cong
\{x,y\}\times\mathcal S_{21}\times\mathcal U_{15}.
}
\] Consequently \[
|T_2(\mathrm{PAB})|=2\cdot21\cdot15=630.
\] The full binary multiplication table is recovered from the owner law
and the two factor transition systems.
\end{theorem}

\begin{proof}
If \(\Theta(f)=\Theta(g)\), then \(f\) and \(g\) agree
on the same-row sector by equality of \(s\), on the normalized cross-row
sector by equality of \(u\), and on the opposite cross-row sector by
\(u^{-}(i)=-u(-i)\). Hence \(\Theta\) is injective.

Surjectivity is T2-REAL: the factorization index lists one realizing
term for every triple in
\(\{x,y\}\times\mathcal S_{21}\times\mathcal U_{15}\). Therefore the
number of binary term functions is \(2\cdot21\cdot15=630\). The product
law follows from \(\operatorname{ow}(f\cdot g)=\operatorname{ow}(f)\)
and the two finite factor transition systems; T2-MULT checks
reconstruction of the full table.
\end{proof}

The factorization separates owner, same-row, and cross-row data. It also
marks the end of the finite term-function layer considered here: the
all-arity term algebra is left for the open-problem list.

The same row/column coordinates used throughout the selector and term
algebra now give a matrix-unit carrier. The next two sections record the
continuous structures supported by that carrier and by the selected
six-edge \(C/J\) geometry.

\section{Carrier-level companion Lie
package}\label{carrier-level-companion-lie-package}

This section records ordinary matrix and Lie structures on the
vector-space carrier exposed after PAB selection
\cite{HornJohnson2012,FultonHarris1991,Hall2015}. Carrier-level means
that the construction uses \[
k[M]\cong\operatorname{Mat}_3(k)
\] and the PAB-compatible simultaneous symbol-permutation action.

\subsection{Carrier identification}\label{carrier-identification}

The carrier vector space identifies with matrix units: \[
k[M]\cong\operatorname{Mat}_3(k),
\qquad
 e_{(r,c)}\leftrightarrow E_{rc}.
\] PAB multiplication and ordinary matrix multiplication are two
bilinear structures on this nine-dimensional carrier. Ordinary matrix
multiplication gives \[
\mathfrak{gl}(3,k)=kI\oplus\mathfrak{sl}(3,k),
\] and the carrier decomposes as \[
k[M]=\mathbf1\oplus\mathbf8
\] under the adjoint matrix Lie package.

\subsection{\texorpdfstring{Compact and split
\(\operatorname{Sym}(S)\)-stable
branches}{Compact and split \textbackslash operatorname\{Sym\}(S)-stable branches}}\label{compact-and-split-operatornamesyms-stable-branches}

Over \(\mathbb R\), the compact three-dimensional branch is \[
\mathfrak{so}(3)=\{X\in\mathfrak{sl}(3,\mathbb R):X^\top=-X\}.
\] It is stable under simultaneous permutation conjugation \[
X\mapsto P_\rho XP_\rho^{-1},
\qquad \rho\in\operatorname{Sym}(S).
\]

There is also a natural split comparator. With \[
\mathbb R^3=\mathbf1\oplus V_{std},
\qquad
V_{std}=\{v:v_0+v_1+v_2=0\},
\] the embedded algebra \[
\mathfrak{sl}(V_{std})\subset\mathfrak{sl}(3,\mathbb R)
\] is the subalgebra with zero row sums and zero column sums. It is
\(\operatorname{Sym}(S)\)-stable and isomorphic to the split real
\(\mathfrak{sl}_2\). Compactness selects the \(\mathfrak{so}(3)\) branch
among these three-dimensional carrier-level branches.

\subsection{Carrier-level compact matrix
package}\label{carrier-level-compact-matrix-package}

\begin{theorem}
\label{thm:19.1}
In the real/complex compact-form setting, the
carrier identification and the PAB-compatible
\(\operatorname{Sym}(S)\)-action support the carrier-level compact
matrix package \[
SU(3)\times SU(2)\times U(1),
\] where \(SU(3)\) is the compact form associated with the
\(\mathfrak{sl}(3)\) matrix package, \(SU(2)\) is the double cover of
the compact \(SO(3)\) branch, and \(U(1)\) is the compact central factor
associated with \(kI\).
\end{theorem}

\begin{proof}
The matrix-unit carrier gives the ordinary matrix
algebra \(\operatorname{Mat}_3\), hence the compact form attached to
\(\mathfrak{sl}(3)\) and the compact central circle attached to \(kI\).
The compact branch \(\mathfrak{so}(3)\) is stable under simultaneous
permutation conjugation, and its double cover supplies the \(SU(2)\)
factor. The bridge audit \texttt{l3\_companion\_symplectic\_summary.csv}
checks the compact type and carrier-level summary.
\end{proof}

\section{Signed symplectic bridge}\label{signed-symplectic-bridge}

This bridge uses the finite pure \(C/J\) directed-edge geometry selected
in Part II. The continuous statement below is a signed symplectic
realization of that six-edge geometry \cite{McDuffSalamon2017}.

Let \(V=\mathbb R^6\) have ordered basis
\[
\mathcal B=
(e_{01},e_{12},e_{20},e_{10},e_{02},e_{21}),
\]
where \(e_{rc}\) denotes the directed edge \((r,c)\). Let
\[
\Omega=
\begin{pmatrix}
0&I_3\\
-I_3&0
\end{pmatrix}
\]
be the matrix of the symplectic form \(\omega\) in this basis. Set
\[
P=
\begin{pmatrix}
0&0&1\\
1&0&0\\
0&1&0
\end{pmatrix},
\qquad
R=
\begin{pmatrix}
1&0&0\\
0&0&1\\
0&1&0
\end{pmatrix}.
\]
The continuation lift, unsigned reversal, sign twist, and signed
reversal are
\[
\widehat C=
\begin{pmatrix}
P&0\\
0&P
\end{pmatrix},
\qquad
\widehat J_0=
\begin{pmatrix}
0&R\\
R&0
\end{pmatrix},
\]
and
\[
S=
\begin{pmatrix}
-I_3&0\\
0&I_3
\end{pmatrix},
\qquad
\widehat J_s=\widehat J_0S=
\begin{pmatrix}
0&R\\
-R&0
\end{pmatrix}.
\]
Here \(\widehat J_0\) is the unsigned reversal on directed edges, while
\(\widehat J_s\) induces the same reversal on projective directed-edge
lines but carries the sign needed for symplecticity.

A \emph{signed projective dihedral lift} of the finite
\[
D_3=\langle C,J\mid C^3=J^2=1,\ JCJ=C^{-1}\rangle
\]
action is a pair of symplectic matrices \(\widehat C,\widehat J_s\in
\operatorname{Sp}(V,\omega)\) whose projective classes define a
homomorphism
\[
D_3\longrightarrow \operatorname{PSp}(V,\omega),
\]
together with an unsigned lift \(\widehat J_0\in GL(V)\) inducing the
same reversal on the six directed-edge lines and satisfying
\[
\widehat J_0^T\Omega\widehat J_0=-\Omega.
\]
Thus the word ``signed'' refers to the replacement of the unsigned
anti-symplectic reversal \(\widehat J_0\) by the signed reversal
\(\widehat J_s\), while ``projective'' records the relation
\(\widehat J_s^2=-I\), which becomes \(J^2=1\) after passing to
\(\operatorname{PSp}(V,\omega)\).

\begin{theorem}
\label{thm:20.1}
The finite pure \(\{C,J\}\) directed-edge geometry admits the signed
projective dihedral lift displayed above. The unsigned reversal is
anti-symplectic, while the signed reversal is symplectic. The finite
tick word
\[
C,J,C,C,J
\]
visits all six directed edges and has a signed symplectic realization.
\end{theorem}

\begin{proof}
Direct matrix multiplication gives
\[
\widehat C^T\Omega\widehat C=\Omega,
\qquad
\widehat J_0^T\Omega\widehat J_0=-\Omega,
\qquad
\widehat J_s^T\Omega\widehat J_s=\Omega.
\]
Moreover
\[
\widehat C^3=I,
\qquad
\widehat J_s^2=-I,
\qquad
\widehat J_s\widehat C\widehat J_s^{-1}=\widehat C^{-1}.
\]
Therefore the projective classes satisfy the dihedral relations in
\(\operatorname{PSp}(V,\omega)\). On the six coordinate lines, the action
of \(\widehat C\) is
\[
(01\ 12\ 20)(10\ 02\ 21),
\]
which is the finite continuation map \(C(r,c)=(c,\overline{rc})\), and
the projective action of \(\widehat J_s\) is the reversal
\[
01\leftrightarrow10,\qquad
12\leftrightarrow21,\qquad
20\leftrightarrow02.
\]
Starting from \(01\), the tick word \(C,J,C,C,J\) gives
\[
01\longmapsto12\longmapsto21\longmapsto10\longmapsto02
\longmapsto20,
\]
so it visits all six directed edges. The bridge audits
\texttt{l3\_signed\_lift\_matrices.csv},
\texttt{l3\_companion\_symplectic\_summary.csv}, and
\texttt{l3\_tick\_map\_audit.csv} record the corresponding finite
checks.
\end{proof}

\part{Certificates and open bridges}

\section{Finite certificate protocol}\label{finite-certificate-protocol}

The support package is available at
\[
\text{\url{https://github.com/burundai-t/sigma-equivariant-magmas-ag23}}.
\]
It separates three roles:
\[
\texttt{scripts/},
\qquad
\texttt{tables/},
\qquad
\texttt{mathcal\_H/}.
\] The compact \texttt{scripts/} and \texttt{tables/} corpus supports
the manuscript-facing Landscape/Selection/Linearization finite checks. The separate
\texttt{mathcal\_H/} bundle supports the global hidden-continuation
theorem of Section 7, and
\texttt{mathcal\_H/h2/verify\_h2\_coverage\_light.py} provides the
static coverage check for the accepted H2 raw artifacts.

The compact check is

\begin{verbatim}
python3 -S scripts/verify_preprint_support.py
\end{verbatim}

It runs inventory and hygiene checks and then V1--V5:

\begin{longtable}[]{@{}
  >{\raggedright\arraybackslash}p{(\columnwidth - 2\tabcolsep) * \real{0.5000}}
  >{\raggedright\arraybackslash}p{(\columnwidth - 2\tabcolsep) * \real{0.5000}}@{}}
\toprule\noalign{}
\begin{minipage}[b]{\linewidth}\raggedright
verifier
\end{minipage} & \begin{minipage}[b]{\linewidth}\raggedright
role
\end{minipage} \\
\midrule\noalign{}
\endhead
\bottomrule\noalign{}
\endlastfoot
V1 & associativity distribution, endpoints, controlled \(\mathcal H\)
comparisons \\
V2 & selection chain, local path-dependence, pure \(C/J\) separation \\
V3 & operator passport, composition envelope, absorption, \(T_2\) \\
V4 & finite bridges: \(C/J\), \(\mathcal H\), matrix, signed
symplectic \\
V5 & auxiliary five-dimensional fiber algebra \\
\end{longtable}

The compact table families are:

\begin{longtable}[]{@{}
  >{\raggedright\arraybackslash}p{(\columnwidth - 4\tabcolsep) * \real{0.3333}}
  >{\raggedright\arraybackslash}p{(\columnwidth - 4\tabcolsep) * \real{0.3333}}
  >{\raggedright\arraybackslash}p{(\columnwidth - 4\tabcolsep) * \real{0.3333}}@{}}
\toprule\noalign{}
\begin{minipage}[b]{\linewidth}\raggedright
family
\end{minipage} & \begin{minipage}[b]{\linewidth}\raggedright
tables
\end{minipage} & \begin{minipage}[b]{\linewidth}\raggedright
purpose
\end{minipage} \\
\midrule\noalign{}
\endhead
\bottomrule\noalign{}
\endlastfoot
L1 landscape atlas & T01--T07 & associativity distribution, endpoints,
controlled \(\mathcal H\) comparison \\
L2 selection & T08--T11 & selection chain and pure \(C/J\) survivor
separation \\
L3 operator/term/bridge & T12--T30 & operator passport, \(T_2\), matrix
bridge, signed symplectic lift, fiber algebra \\
\end{longtable}

The complete file-by-file index, including H2 light coverage support, is
\texttt{SUPPORT\_INDEX.txt}.

\subsection{Global hidden-continuation certificate handles (Stage 6
bundle)}\label{global-hidden-continuation-certificate-handles-stage-6-bundle}

The global hidden-continuation theorem is checked by the separate bundle

\begin{verbatim}
cd mathcal_H
PYTHONDONTWRITEBYTECODE=1 python3 verify_stage6_bundle.py
\end{verbatim}

The manuscript-level implication is \[
\text{local distance reduction}
+
\text{H1/H2/H3/H4 certificates}
\Longrightarrow
\text{Theorem 7.2}.
\]

\begin{longtable}[]{@{}
  >{\raggedright\arraybackslash}p{(\columnwidth - 4\tabcolsep) * \real{0.3333}}
  >{\raggedright\arraybackslash}p{(\columnwidth - 4\tabcolsep) * \real{0.3333}}
  >{\raggedright\arraybackslash}p{(\columnwidth - 4\tabcolsep) * \real{0.3333}}@{}}
\toprule\noalign{}
\begin{minipage}[b]{\linewidth}\raggedright
handle
\end{minipage} & \begin{minipage}[b]{\linewidth}\raggedright
bundle location
\end{minipage} & \begin{minipage}[b]{\linewidth}\raggedright
theorem component
\end{minipage} \\
\midrule\noalign{}
\endhead
\bottomrule\noalign{}
\endlastfoot
S6-RED, S6-CERT-IF & \texttt{stage6.md}, verifier & local distance
reduction \\
S6-H1 & \texttt{h1/} & unrestricted full-landscape range \\
S6-H2 & \texttt{h2/} & pure-frontier UNSAT certificate \\
S6-H3 & embedded verifier check & PAB/row-complement pure-frontier
witnesses \\
S6-H4 & \texttt{h4/} & signed-cancellation spike classifier \\
\end{longtable}

Focused checks are:

\begin{verbatim}
PYTHONDONTWRITEBYTECODE=1 python3 verify_stage6_bundle.py --only-h1
PYTHONDONTWRITEBYTECODE=1 python3 verify_stage6_bundle.py --only-certpack-final
PYTHONDONTWRITEBYTECODE=1 python3 verify_stage6_bundle.py --only-h3
PYTHONDONTWRITEBYTECODE=1 python3 verify_stage6_bundle.py --only-h4
\end{verbatim}

The H2 obligation proves \[
(h_q\ge0\ \forall q\in S^5)\wedge rawH\ge1171
\] UNSAT on the depth-9 live frontier: \[
18540/18540\ \mathrm{UNSAT},
\qquad SAT=0,
\qquad UNKNOWN=0.
\] The light static H2 coverage verifier checks the accepted raw H2
artifact coverage and integrity for the declared live frontier:

\begin{verbatim}
cd mathcal_H/h2
PYTHONDONTWRITEBYTECODE=1 python3 verify_h2_coverage_light.py
\end{verbatim}

\subsection{Dependency summary}\label{dependency-summary}

\begin{longtable}[]{@{}
  >{\raggedright\arraybackslash}p{(\columnwidth - 2\tabcolsep) * \real{0.5000}}
  >{\raggedright\arraybackslash}p{(\columnwidth - 2\tabcolsep) * \real{0.5000}}@{}}
\toprule\noalign{}
\begin{minipage}[b]{\linewidth}\raggedright
result family
\end{minipage} & \begin{minipage}[b]{\linewidth}\raggedright
support role
\end{minipage} \\
\midrule\noalign{}
\endhead
\bottomrule\noalign{}
\endlastfoot
compact PAB uniqueness, Theorem 12.3 & hand finite proof; L2 audit
tables \\
explanatory selection chain, Proposition 12.4 & didactic chain,
logically subsumed by Theorem 12.3; local audits \\
global associativity atlas & finite atlas tables T01--T05 \\
global hidden-continuation frontier & separate H1--H4 certificate
bundle \\
operator passport & analytic/classification results with L3 audits \\
binary term algebra & sector reduction plus realization/transition
certificates \\
matrix/Lie and signed symplectic structures & bridge audits \\
\end{longtable}

\section{Open problems and bridges}\label{open-problems-and-bridges}

The following fronts are subsequent work. The selection theorem and the
operator passport above are complete at the level stated in this
manuscript.

\subsection{Closures}\label{closures}

\subsubsection{\texorpdfstring{Functorial PAB--\(\mathfrak{gl}(3)\)
bridge}{Functorial PAB--\textbackslash mathfrak\{gl\}(3) bridge}}\label{functorial-pabmathfrakgl3-bridge}

The carrier identification \[
k[M]\cong\operatorname{Mat}_3(k)
\] supports two bilinear operations: PAB multiplication and ordinary
matrix multiplication. The cross-row contraction \[
\mu_{cross}(A,B)=AJ_3B^\top
\] and the carrier-level compact package give partial alignments. The
closure question is whether there is a universal or functorial
construction from the PAB magma to the matrix/Lie package, beyond the
shared carrier.

\subsubsection{\texorpdfstring{Symplectic development of pure
\(C/J\)}{Symplectic development of pure C/J}}\label{symplectic-development-of-pure-cj}

The finite pure \(C/J\) geometry has an explicit signed symplectic lift.
The next problem is to develop the associated continuous symplectic
geometry, including canonical \((q,p)\)-charts, Hamiltonian paths, and
coordinate-free separation of pure reversal from drifted reversal.

\subsubsection{Unified selection principle}\label{unified-selection-principle}

Theorem~\ref{thm:12.3} selects PAB by the conjunction of four finite
selectors of distinct nature: information-theoretic, entropic,
algebraic, and combinatorial-geometric. It remains open whether these
selectors arise from a single natural, non-tautological
information-theoretic or algebraic functional whose global minimizers on
\(\Omega'\) are exactly the PAB representative. In such a formulation,
PAB would appear as a global extremum of one canonical object rather
than as the intersection of four independently stated extrema.

\subsection{Extensions}\label{extensions}

\subsubsection{All-arity term algebra}\label{all-arity-term-algebra}

The binary term algebra satisfies \[
|T_2(\mathrm{PAB})|=630,
\qquad
T_2(\mathrm{PAB})\cong\{x,y\}\times\mathcal S_{21}\times\mathcal U_{15}.
\] The all-arity term algebra remains open: the main question is whether
the binary factorization extends to a general clone structure or is
genuinely binary \cite{Post1941,Szendrei1986,Lau2006}.

\subsubsection{Larger carriers}\label{larger-carriers}

Natural generalizations include analogues over larger carriers
\(S\times S\), especially finite fields or cyclic sets, and
directed-edge \(C/J\) geometries beyond the triangle.

\subsubsection{Removing AX2}\label{removing-ax2}

The landscape \(\Omega'\) fixes the same-row Steiner axiom. The larger
landscape without AX2 asks whether the same-row Steiner rule can itself
be selected by intrinsic criteria.

\subsection{Status summary}\label{status-summary}

\begin{longtable}[]{@{}
  >{\raggedright\arraybackslash}p{(\columnwidth - 2\tabcolsep) * \real{0.5000}}
  >{\raggedright\arraybackslash}p{(\columnwidth - 2\tabcolsep) * \real{0.5000}}@{}}
\toprule\noalign{}
\begin{minipage}[b]{\linewidth}\raggedright
Front
\end{minipage} & \begin{minipage}[b]{\linewidth}\raggedright
Status
\end{minipage} \\
\midrule\noalign{}
\endhead
\bottomrule\noalign{}
\endlastfoot
Global hidden-continuation certificate & closed by Theorem 7.2 and the
\texttt{mathcal\_H/} certificate bundle \\
Functorial PAB--\(\mathfrak{gl}(3)\) bridge & open \\
Unified selection principle & open \\
Symplectic development of pure \(C/J\) & open \\
All-arity term algebra & open \\
Larger carriers & open \\
Removing AX2 & open \\
\end{longtable}

\section{Conclusion}\label{conclusion}

The PAB magma is selected from the finite normal-form landscape
\(\Omega'\cong S^{21}\) by four internal finite criteria: \[
\mathcal H_{acc}=0,
\qquad
H_{diag}=0,
\qquad
\text{diagonal idempotence},
\qquad
\text{pure }C/J.
\] The surrounding landscape is mapped by certified associativity and
hidden-continuation atlases, and the selected magma then exposes its
linear/operator passport, binary term algebra, carrier-level compact Lie
package, and signed symplectic bridge.

Thus the theorem package has the intended three-layer form: \[
\boxed{\text{finite landscape}}
\qquad
\boxed{\text{intrinsic finite selection}}
\qquad
\boxed{\text{linearization of PAB}}.
\]

\section*{Note on AI assistance}
\addcontentsline{toc}{section}{Note on AI assistance}

This work was developed with the assistance of large language models.

The author specified the conceptual framework, constraints, and target
properties of the objects under study. Within this framework, AI tools
were used for computational exploration, code drafting, and structuring
intermediate results.

All mathematical definitions, statements, and theorems presented in the
paper were established and independently verified by the author. The
finite claims are backed by exhaustive checks recorded in the
accompanying support package.

\newpage
\begin{thebibliography}{99}
\bibitem{AppelHaken1977} K. Appel and W. Haken, ``Every planar map is four colorable. Part I:
Discharging,'' \emph{Illinois Journal of Mathematics} \textbf{21}
(1977), no. 3, 429--490. doi:10.1215/ijm/1256049011.

\bibitem{AppelHakenKoch1977} K. Appel, W. Haken, and J. Koch, ``Every planar map is four colorable.
Part II: Reducibility,'' \emph{Illinois Journal of Mathematics}
\textbf{21} (1977), no. 3, 491--567. doi:10.1215/ijm/1256049012.

\bibitem{Bennett1973} C. H. Bennett, ``Logical reversibility of computation,'' \emph{IBM
Journal of Research and Development} \textbf{17} (1973), no. 6,
525--532. doi:10.1147/rd.176.0525.

\bibitem{Bennett1982} C. H. Bennett, ``The thermodynamics of computation---a review,''
\emph{International Journal of Theoretical Physics} \textbf{21} (1982),
no. 12, 905--940. doi:10.1007/BF02084158.

\bibitem{Birkhoff1935} G. Birkhoff, ``On the structure of abstract algebras,''
\emph{Proceedings of the Cambridge Philosophical Society} \textbf{31}
(1935), no. 4, 433--454. doi:10.1017/S0305004100013463.

\bibitem{Bruck1958} R. H. Bruck, \emph{A Survey of Binary Systems}, Ergebnisse der
Mathematik und ihrer Grenzgebiete, Neue Folge, Heft 20, Springer,
Berlin, 1958.

\bibitem{BurrisSankappanavar1981} S. Burris and H. P. Sankappanavar, \emph{A Course in Universal Algebra},
Graduate Texts in Mathematics 78, Springer-Verlag, New York, 1981.
doi:10.1007/978-1-4613-8130-3.

\bibitem{Cohn1965} P. M. Cohn, \emph{Universal Algebra}, Harper's Series in Modern
Mathematics, Harper \& Row, New York, 1965.

\bibitem{ColbournDinitz2007} C. J. Colbourn and J. H. Dinitz, eds., \emph{Handbook of Combinatorial
Designs}, 2nd ed., Chapman \& Hall/CRC, Boca Raton, 2007.

\bibitem{CoverThomas2006} T. M. Cover and J. A. Thomas, \emph{Elements of Information Theory}, 2nd
ed., Wiley-Interscience, Hoboken, 2006. doi:10.1002/047174882X.

\bibitem{DenesKeedwell1974} J. Dénes and A. D. Keedwell, \emph{Latin Squares and Their
Applications}, Academic Press, New York, 1974.

\bibitem{FultonHarris1991} W. Fulton and J. Harris, \emph{Representation Theory: A First Course},
Graduate Texts in Mathematics 129, Springer-Verlag, New York, 1991.
doi:10.1007/978-1-4612-0979-9.

\bibitem{Gonthier2008} G. Gonthier, ``Formal proof---the four-color theorem,'' \emph{Notices of
the American Mathematical Society} \textbf{55} (2008), no. 11,
1382--1393.

\bibitem{GuilleminSternberg1990} V. Guillemin and S. Sternberg, \emph{Symplectic Techniques in Physics},
Cambridge University Press, Cambridge, 1990.

\bibitem{Hales2005} T. C. Hales, ``A proof of the Kepler conjecture,'' \emph{Annals of
Mathematics} \textbf{162} (2005), no. 3, 1065--1185.
doi:10.4007/annals.2005.162.1065.

\bibitem{HalesEtAl2017} T. Hales, M. Adams, G. Bauer, T. D. Dang, J. Harrison, L. T. Hoang, C.
Kaliszyk, V. Magron, S. McLaughlin, T. T. Nguyen, Q. T. Nguyen, T.
Nipkow, S. Obua, J. Pleso, J. Rute, A. Solovyev, T. H. A. Ta, N. T.
Tran, T. D. Trieu, J. Urban, K. Vu, and R. Zumkeller, ``A formal proof
of the Kepler conjecture,'' \emph{Forum of Mathematics, Pi} \textbf{5}
(2017), e2. doi:10.1017/fmp.2017.1.

\bibitem{Hall2015} B. C. Hall, \emph{Lie Groups, Lie Algebras, and Representations: An
Elementary Introduction}, 2nd ed., Graduate Texts in Mathematics 222,
Springer, Cham, 2015. doi:10.1007/978-3-319-13467-3.

\bibitem{HornJohnson2012} R. A. Horn and C. R. Johnson, \emph{Matrix Analysis}, 2nd ed., Cambridge
University Press, Cambridge, 2012.

\bibitem{Landauer1961} R. Landauer, ``Irreversibility and heat generation in the computing
process,'' \emph{IBM Journal of Research and Development} \textbf{5}
(1961), no. 3, 183--191. doi:10.1147/rd.53.0183.

\bibitem{Lau2006} D. Lau, \emph{Function Algebras on Finite Sets: A Basic Course on
Many-Valued Logic and Clone Theory}, Springer Monographs in Mathematics,
Springer, Berlin, 2006. doi:10.1007/3-540-36023-9.

\bibitem{McCune1997} W. McCune, ``Solution of the Robbins problem,'' \emph{Journal of
Automated Reasoning} \textbf{19} (1997), no. 3, 263--276.
doi:10.1023/A:1005843212881.

\bibitem{McDuffSalamon2017} D. McDuff and D. Salamon, \emph{Introduction to Symplectic Topology},
3rd ed., Oxford Graduate Texts in Mathematics 27, Oxford University
Press, Oxford, 2017.

\bibitem{Pflugfelder1990} H. O. Pflugfelder, \emph{Quasigroups and Loops: Introduction}, Sigma
Series in Pure Mathematics 7, Heldermann Verlag, Berlin, 1990.

\bibitem{Post1941} E. L. Post, \emph{The Two-Valued Iterative Systems of Mathematical
Logic}, Annals of Mathematics Studies 5, Princeton University Press,
Princeton, 1941.

\bibitem{Rosenberg1970} I. G. Rosenberg, ``Über die funktionale Vollständigkeit in den
mehrwertigen Logiken,'' \emph{Rozpravy Československé Akademie Věd, Řada
matematických a přírodních věd} \textbf{80} (1970), no. 4, 3--93.

\bibitem{Shannon1948} C. E. Shannon, ``A mathematical theory of communication,'' \emph{Bell
System Technical Journal} \textbf{27} (1948), 379--423 and 623--656.
doi:10.1002/j.1538-7305.1948.tb01338.x;
doi:10.1002/j.1538-7305.1948.tb00917.x.

\bibitem{Stanley2012} R. P. Stanley, \emph{Enumerative Combinatorics, Volume 1}, 2nd ed.,
Cambridge Studies in Advanced Mathematics 49, Cambridge University
Press, Cambridge, 2012. doi:10.1017/CBO9781139058520.

\bibitem{Szendrei1986} Á. Szendrei, \emph{Clones in Universal Algebra}, Séminaire de
mathématiques supérieures 99, Presses de l'Université de Montréal,
Montréal, 1986.
\end{thebibliography}

\end{document}
